% SPDX-FileCopyrightText: 2026 Will Cook
% SPDX-License-Identifier: CC-BY-4.0
%
% Standalone Erdős Problem Note for #257.  Context is deliberately repeated
% from the gateway so that this file remains independently readable.
%
% Ordering rule, revised July 2026.  The previous draft was a map of settled
% supports with a method boundary attached.  A map is a list, and #257 is a
% universal statement; adding rows does not approach one.  This draft leads
% with the finite theorem, then with the statements that quantify over *every*
% hypothetical rational-valued support, and demotes the settled families to
% context.  Those forced-structure theorems were already in the corpus and
% were simply absent from the note.
%
% Second correction.  The squarefree support was printed here as open.  It is
% known at every power-of-two base: Duverney and Tachiya settled the joint
% family in 2019, and this note
% had failed to cite them.  The section that claimed a frontier now records
% what the position actually is -- a certificate engine failing on a value
% that is known irrational, for a reason internal to a normalisation
% convention rather than to the value.  That is a smaller claim and a more
% useful one.
\documentclass[11pt]{article}
\input{problem-note-preamble}
\renewcommand{\commit}{5e28e73cf70a0ccf3219bb7197401307fb6d3bfd}
\renewcommand{\commitshort}{5e28e73cf70a}
\hypersetup{
  pdftitle={Reciprocal-summable support irrationality at every integer base (Erdős Problem 257)},
  pdfsubject={Erdős Problem Notes: Problem 257},
  pdfkeywords={irrationality, Mersenne numbers, Lambert series, divisor incidence, scaled tails, achievement sets, squarefree, Lean 4},
}

\newcommand{\wgt}[1]{w_{#1}}
\newcommand{\Ach}{\mathcal A}
\newcommand{\scA}{\operatorname{sc}}
\newcommand{\Asf}{A_{\mathrm{sf}}}

\runninghead{Erd\H{o}s \#257: all-base reciprocal-summable supports}
\title{Reciprocal-Summable Support Irrationality at Every Integer Base}
\subtitle{The series $\sum_{a\in A}(b^a-1)^{-1}$ and Erd\H{o}s Problem~\#257}
\author{Will Cook}
\date{July 2026}

\begin{document}
\maketitle
\noteprovenance

\begin{abstract}
\noindent
Erd\H{o}s Problem~\#257 asks whether
\[
 X_A=\sum_{a\in A}\frac1{2^a-1}
\]
is irrational for every infinite \(A\subseteq\Npos\).  The problem remains
open.  The leading Lean-checked support-family theorem covers every infinite
\(A\subseteq\Npos\) with convergent reciprocal mass, simultaneously at every
integer base \(b\ge2\):
\[
 \sum_{a\in A}\frac1a<\infty
 \quad\Longrightarrow\quad
 \sum_{a\in A}\frac1{b^a-1}\notin\Q
 \qquad(b\ge2).
\]
No pairwise-coprimality, periodicity or density hypothesis is imposed.  The
proof first obtains arbitrarily close positive binary returns of the shifted
divisor-tail sum from exact gcd-orbit means, an LCM-prefix limit and Cesaro
averaging.  A uniform atom inequality transfers every such return to every
larger radix, and an exact radix-\(b\) integral scaled-tail gap rules out
rationality.  Thus every counterexample to Problem~\#257 must have divergent
reciprocal mass.  In particular every infinite set of powerful integers is
included, since the representation \(n=x^2y^3\) is dominated by
\(\zeta(2)\zeta(3)\).  Erd\H{o}s printed this theorem in 1968 for pairwise-coprime
support at every integer base \(t\ge2\), and said in the next sentence that
coprimality can be removed without printing that argument.  The all-base
quantifier is therefore his.  The coprimality-free conclusion is the extension
he stated; reciprocal summability remains the hypothesis.  No identification
with Erd\H{o}s's omitted argument and no theorem-priority claim are made.

For a finite nonempty support \(F\), the reduced denominator \(D_F\) of
\(\sum_{n\in F}(b^n-1)^{-1}\) satisfies
\[
 \operatorname{ord}_{D_F}(b)=\operatorname{lcm}(F);
\]
if \(\operatorname{lcm}(F)\ge2\), then
\(\operatorname{lcm}(F)<D_F\).  For a hypothetical rational value on an
infinite support, binary long division produces an integral scaled-tail orbit
obeying an exact recurrence.  Unbounded shifted tails and uniformly bounded
zero windows hold for every infinite or nonempty support respectively, with no
rationality hypothesis; they are not advertised here as restrictions on
rational counterexamples.  Boolean--M\"obius inversion gives an
exact orbit characterisation of all rational support values, but does not make
the reconstructed support infinite.

The achievement set of the weights \((2^n-1)^{-1}\) is compact, perfect,
totally disconnected, nowhere dense, and of Lebesgue measure \(1\).  If the
allowed exponent set has finite complement \(F\), the restricted achievement
set has measure \(2^{-|F|}\); if the complement is infinite, it has measure
zero.  These geometric results classify sets of subsums, not the arithmetic
nature of their points.

At the counterexample frontier, a sequence of finite Boolean support words
whose depths tend to infinity and whose absolute terminal carry is
\(o(2^M)\) produces an infinite support with value exactly \(1/2\), thereby
refuting the universal irrationality assertion.  This is a sharp conditional
endpoint: no such terminal scaled-vanishing sequence is constructed here.

A second conditional endpoint retains the actual integer greedy dynamics.
For every realized upper transition and its literal right run, the complete
base-four affine packet lower bound is equivalent to one pulse-free lower
bound on the immediate successor remainder.  Either universal form produces
an infinite support of value exactly \(1/2\) and hence the same refutation.
No such lower-envelope producer is proved here; this result exposes its exact
successor boundary rather than supplying a counterexample.

Two rational targets make the remaining difficulty explicit.  Membership of
\(1/2\) is equivalent to infinitely many omissions in its greedy expansion;
non-membership is witnessed by one finite fatal gap.  Membership of \(1/21\)
is equivalent to cofinal crossings of the exact scaled lower separatrix, or
equivalently to failure of the stated fatal aligned branch.  Neither
membership question is resolved.  A representation of either value must have
infinite support.

\end{abstract}

\section{Introduction and main results}\label{sec:problem}

\begin{problem}[Erd\H{o}s \#257]\label{res:problem}
For every infinite $A\subseteq\Npos$, is
$X_A=\sum_{a\in A}\dfrac1{2^a-1}$ irrational?
\end{problem}

\noindent See Erd\H{o}s~\cite[p.~222]{erdos1968}, Erd\H{o}s and
Graham~\cite[p.~62]{erdosgraham1980}, and
Erd\H{o}s~\cite[p.~105]{erdos1988}.  Bloom's current catalogue record
reproduces the displayed problem and labels it open, while explicitly warning
that the status is the website owner's present assessment and may omit
relevant literature~\cite{erdosproblems}.  We therefore use the catalogue for
numbering and current reported status only; the original publications and the
later cited papers carry the mathematical claims.

Write
$\wgt n=(2^n-1)^{-1}$.  Expanding each weight as a geometric series and
interchanging the two nonnegative sums gives the coordinate this note works
in.  With the \emph{divisor incidence}
$\scA_A(n)=\#\{d\mid n: d\in A\}$,
\[
  X_A=\sum_{a\in A}\frac1{2^a-1}=\sum_{n\ge1}\frac{\scA_A(n)}{2^n}.
  \tag{1}\label{eq:incidence}
\]
The transform is worth reading carefully, because it is where the arithmetic
of the problem enters.  The datum is a $0/1$ selector, the indicator of $A$;
what appears in~\eqref{eq:incidence} is not that selector but its divisor
transform, a nonnegative integer sequence bounded by the divisor function,
$\scA_A(n)\le d(n)$.  So \#257 is not a generic question about binary digit
sequences.  Equation~\eqref{eq:incidence} is a power-series representation,
not a binary expansion: its coefficients are divisor counts drawn from a
single support and may exceed~$1$.  Every theorem below is a statement about
sequences of that shape.

A small instance fixes the notation.  At $A=\{2,3\}$ the incidence sequence
begins
\[
  \scA_A(1),\dots,\scA_A(12)=0,1,1,1,0,2,0,1,1,1,0,2,
\]
the value $2$ occurring at the multiples of $6$ and the value $0$ at the
integers prime to $6$, and~\eqref{eq:incidence} reads
$\tfrac13+\tfrac17=\tfrac{10}{21}$.

Several statements below hold at every integer base, so we write
$X_A(b)=\sum_{a\in A}(b^a-1)^{-1}$ for the value at an integer base $b\ge2$,
with $X_A=X_A(2)$.  Base~$2$ is the case Problem~\ref{res:problem} asks about
and is meant whenever no base is named.

\paragraph{The reciprocal-summable support theorem.}
This is the result closest to the universal quantifier in
Problem~\ref{res:problem} that is proved here.

\begin{theorem}[reciprocal-summable supports]\label{res:reciprocal-support}
Let $A\subseteq\Npos$ be infinite.  If
\[
  \sum_{a\in A}\frac1a<\infty,
\]
then \(X_A(b)\) is irrational for every integer \(b\ge2\).
\end{theorem}

\noindent The exact checked declaration is
\begin{quote}\small
\texttt{Erdos257PeriodNoncollapse.}\\
\texttt{irrational\_erdosSupportSeries\_of\_summable\_reciprocal}.
\end{quote}
Here the Lean hypothesis is
\texttt{Summable (reciprocalSupportTerm A)}; because $A\subseteq\Npos$,
this is precisely convergence of the displayed nonnegative reciprocal sum.
The formal theorem has no pairwise-coprimality, periodicity or density
assumption.  Its all-base quantifier is part of the checked statement, not a
corollary asserted only in prose.

\paragraph{The finite-support theorem.}
For a finite nonempty $F\subseteq\Npos$ and an integer $b\ge2$, write
\[
  x_F(b)=\sum_{n\in F}\frac1{b^n-1}=\frac{N_F}{D_F}
  \qquad (\gcd(N_F,D_F)=1,\ D_F>0).
\]
We use the standard trivial-modulus convention
$\operatorname{ord}_1(b)=1$.

\begin{theorem}[finite-period noncollapse]\label{res:period}
Let $F\subseteq\Npos$ be finite and nonempty, let $b\ge2$ be an integer, and
let $D_F>0$ be the denominator of $x_F(b)$ in lowest terms.  Then $D_F$ is
coprime to $b$, and
\[
  \operatorname{ord}_{D_F}(b)=\lcm\{n:n\in F\}.
\]
If moreover $\lcm(F)\ge2$, then $\lcm(F)<D_F$.
\end{theorem}

\noindent The order statement is
\mword{Erdos249257/CertificateKernel.lean}{5091}{finite_period_noncollapse}{checked},
its reduced-denominator form is
\mword{Erdos249257/CertificateKernel.lean}{5246}{finite_period_noncollapse_rat_den}{checked},
coprimality is
\mword{Erdos249257/CertificateKernel.lean}{5221}{coprime_base_den_finiteErdosSum}{checked},
and the growth clause is
\mword{Erdos249257/CertificateKernel.lean}{5260}{lcm_lt_den_finiteErdosSum}{checked}.

\paragraph{The conditional half-value endpoint.}
Theorem~\ref{res:terminalhalf} proves that terminal scaled vanishing, meaning finite
Boolean support words at depths \(M\to\infty\), with ranks zero and one
absent and terminal carry \(o(2^M)\), already yields an infinite support
whose reciprocal-Mersenne sum is exactly \(1/2\).  Hence that hypothesis
refutes the universal assertion in Problem~\ref{res:problem}.  This is the
weakest sequential terminal-only producer currently used by the paper, but
its existence remains open; the theorem is a conditional counterexample
endpoint, not a counterexample.

\paragraph{The actual upper-successor endpoint.}
Theorem~\ref{res:actual-upper-successor} keeps the concrete integer greedy
remainder and every pulse of a realized upper/right run.  It proves that the
complete terminal packet inequality is equivalent to the immediate-successor
bound
\[
  2^{d+1}-2^{d-k+1}+2(d+k)\le R_{d+1}.
\]
Thus the accumulated pulse register can be eliminated without weakening this
producer.  If the bound holds for every mature realized run, an infinite
support sums to \(1/2\) and Problem~\ref{res:problem}'s universal assertion is
false.  The universal bound itself remains open, so this is a conditional
counterexample endpoint and an exact target for the remaining Lean search.

\statusboundary

\paragraph{Structure.}
Section~\ref{sec:reciprocal-support} gives the proof architecture of
Theorem~\ref{res:reciprocal-support}, including its exact integer-gap
endgame and the binary-to-radix transfer.  Section~\ref{sec:period} develops the arithmetic of
Theorem~\ref{res:period}.  Section~\ref{sec:forced} collects what
rationality would force on an arbitrary infinite support, which is the part of
this note that quantifies over every support rather than sampling.
Section~\ref{sec:map} gives representative supports already known to give
irrational values, grouped by the argument that reaches them.
Section~\ref{sec:squarefree} treats
the squarefree support, whose values are known at every power-of-two base and which two of the
arguments used here provably cannot reach at any even base.
Section~\ref{sec:geometry} gives the unrestricted and support-restricted
topology and measure classifications, followed by the exact $1/2$ and $1/21$
membership criteria and unresolved branches.  Section~\ref{sec:open} states
what remains.

\section{Reciprocal-summable supports at every integer base}
\label{sec:reciprocal-support}

The hard step in Theorem~\ref{res:reciprocal-support} is not a sparse
denominator estimate.  It is a recurrence argument which makes the initial
radix-\(b\) divisor tail a strict minimum, together with a uniform comparison
which transfers arbitrarily close binary returns to every larger radix.

Let \(c_A(n)=\#\{d\mid n:d\in A\}\).  For \(b\ge2\), put
\[
 T_N^{(b)}=\sum_{n\ge1}\frac{c_A(N+n)}{b^n},
 \qquad
 \omega_{b,d}(N)=\frac{b^{N\bmod d}}{b^d-1}.
\]
Expanding each selected denominator geometrically and regrouping the
nonnegative terms gives the shifted-atom identity
\[
 T_N^{(b)}=\sum_{d\in A}\omega_{b,d}(N).
 \tag{2}\label{eq:shifted-support-atoms}
\]
In particular \(T_0^{(b)}=X_A(b)\).  Every atom is minimized at shift zero.
If \(A\) is infinite and \(N>0\), choose \(d\in A\) with \(d>N\); then the
\(d\)-atom is strictly larger at \(N\), and hence
\[
 T_0^{(b)}<T_N^{(b)}. \tag{3}\label{eq:strict-initial-minimum}
\]

We first work at \(b=2\).  For a positive sampling step \(Q\), the exact mean
of the \(d\)-th binary atom along one orbit of positive \(Q\)-multiples is
\[
 M_Q(d)=\frac{\gcd(Q,d)}{d\bigl(2^{\gcd(Q,d)}-1\bigr)}.
 \tag{4}\label{eq:gcd-orbit-mean}
\]
It satisfies \(0\le M_Q(d)\le1/d\).  Take
\(Q_P=\operatorname{lcm}(1,\ldots,P)\).  For each fixed \(d\), eventually
\(d\mid Q_P\), so \(M_{Q_P}(d)\) tends to \((2^d-1)^{-1}\).
The reciprocal-mass hypothesis is exactly the summable majorant needed to
pass this limit through the sum over \(d\in A\):
\[
 \sum_{d\in A}M_{Q_P}(d)\longrightarrow T_0^{(2)}.
 \tag{5}\label{eq:lcm-prefix-orbit-limit}
\]
For fixed \(Q>0\), the atomwise Cesaro means are dominated by the summable
comparison
\[
 \frac1X\sum_{m=1}^X\omega_{2,d}(Qm)
 =\sum_{r\ge1}2^{-r}\frac{\#\{m\le X:d\mid Qm+r\}}X
 \le\sum_{r\ge1}2^{-r}\frac{Q+r}{d}=\frac{Q+2}{d}.
\]
The numerator counts distinct positive multiples of \(d\) no larger than
\(QX+r\).  Reciprocal summability of \(A\) therefore licenses passage from
atomwise orbit averages to the Cesaro means of the complete shifted sums
\(T_{(k+1)Q}^{(2)}\), which tend to \(\sum_{d\in A}M_Q(d)\).  This domination is
expository; the close-return theorem is already Lean-checked and is not a new
kernel obligation.  Choosing \(P\), then a sufficiently long finite average,
and finally one term no larger than that average proves
\[
 \forall\varepsilon>0\ \exists N>0:\qquad
  T_0^{(2)}<T_N^{(2)}<T_0^{(2)}+\varepsilon.
 \tag{6}\label{eq:strict-close-return}
\]

The all-base step is a pointwise comparison, not a repetition of the orbit
calculation.  If \(r=N\bmod d\), direct subtraction gives
\[
 \omega_{b,d}(N)-\omega_{b,d}(0)
   =\frac{b^r-1}{b^d-1}.
\]
For \(r=0\) both sides below vanish.  For \(0<r<d\),
\[
 \frac{b^r-1}{b^d-1}
 \le b^{-(d-r)}
 \le 2^{-(d-r)}
 \le 2\,\frac{2^r-1}{2^d-1}.
\]
Consequently, uniformly in \(b\ge2\), \(N\), and \(d\),
\[
 0\le T_N^{(b)}-T_0^{(b)}
 \le2\bigl(T_N^{(2)}-T_0^{(2)}\bigr).
 \tag{7}\label{eq:binary-radix-transfer}
\]
Applying the binary close-return statement with tolerance
\(\varepsilon/2\), then using~\eqref{eq:strict-initial-minimum}, yields
\[
 \forall b\ge2\ \forall\varepsilon>0\ \exists N>0:\qquad
 T_0^{(b)}<T_N^{(b)}<T_0^{(b)}+\varepsilon.
 \tag{8}\label{eq:all-base-close-return}
\]
The Lean-checked declarations corresponding to the atom comparison and the
transferred close return are
\mword{Erdos249257/AllBaseReciprocalSupportIrrationality.lean}{270}{shiftedRadixAtom\_sub\_zero\_le\_two\_mul\_binary\_sub\_zero}{Lean}
and
\mword{Erdos249257/AllBaseReciprocalSupportIrrationality.lean}{362}{exists\_shiftedRadixSupportAtom\_closeReturn\_of\_summable\_reciprocal}{Lean}.

It remains to turn a close return into an arithmetic contradiction.  The atom
recurrence sums exactly to
\[
 bT_N^{(b)}-T_{N+1}^{(b)}=c_A(N+1).
 \tag{9}\label{eq:radix-tail-recurrence}
\]
Suppose \(T_0^{(b)}=X_A(b)=p/v\) with \(v>0\).  Starting from \(z_0=p\), define
\[
 z_{N+1}=bz_N-vc_A(N+1).
\]
Equation~\eqref{eq:radix-tail-recurrence} gives
\(z_N=vT_N^{(b)}\in\mathbb Z\) for every \(N\).  Apply
\eqref{eq:all-base-close-return} with \(\varepsilon=1/v\).  Then
\[
 0<z_N-z_0=v\bigl(T_N^{(b)}-T_0^{(b)}\bigr)<1,
\]
which is impossible for an integer.  This proves
Theorem~\ref{res:reciprocal-support}.  The exact summed recurrence and final
theorem are
\mword{Erdos249257/AllBaseReciprocalSupportIrrationality.lean}{200}{tsum\_shiftedRadixSupportAtom\_step}{Lean}
and
\mword{Erdos249257/AllBaseReciprocalSupportIrrationality.lean}{395}{irrational\_erdosSupportSeries\_of\_summable\_reciprocal}{Lean}.
Comparator exposes only the final irrationality theorem; the transfer,
close-return, and integer-orbit results remain its load-bearing subordinate
mechanism rather than theorem-count padding.

One immediate consequence is worth making explicit.  The reciprocal sum over
all powerful integers converges (use the standard representation
\(n=x^2y^3\), with \(y\) squarefree, and dominate by
\(\zeta(2)\zeta(3)\)).  Therefore every infinite support contained in the
powerful integers, with no isolated-prime-power hypothesis at all, has
irrational \(X_A(b)\) at every integer base \(b\ge2\).  This is a consequence
of the general theorem, not a separate Comparator endpoint.

\paragraph{Attribution and boundary.}\label{bdry:reciprocal-support}
Erd\H{o}s's printed theorem on p.~222 reads: if \((n_i,n_j)=1\) and
\(\sum1/n_i<\infty\), then \(\sum1/(t^{n_i}-1)\) is irrational for every
\(t\ge2\).  It is already an all-base theorem.  The next sentence says that
the condition \((n_i,n_j)=1\) is superfluous and that the details are not
given; p.~226 repeats that boundary.  The single hypothesis
Theorem~\ref{res:reciprocal-support} removes relative to the printed theorem
is pairwise coprimality, which is exactly the hypothesis Erd\H{o}s said he
could remove.  The base generality is his, not ours.  What is added is a
complete machine-checked proof of the withheld case.  We do not identify this
proof with Erd\H{o}s's omitted proof or make a theorem-priority claim.
The theorem says nothing about supports with divergent reciprocal mass.
Consequently it leaves Problem~\ref{res:problem} open, but it removes the
entire reciprocal-summable region from the possible counterexample space.

\section{Extensions beyond reciprocal summability}
\label{sec:eight-return-extensions}

The following arguments are ordinary mathematical proofs.  Their complete
analytic formalizations are not part of the checked theorem above.  They
exclude several further classes with divergent reciprocal mass; they do not
decide the universal problem.  Complete proofs, examples and countermodels
are retained with this paper in
\href{https://github.com/wcook04/plectis-lean-erdos249-257/blob/main/docs/research/erdos257-eight-return-proofs.md}{the analytic proof supplement}.
The variable-exponent argument below and the finite repair deadline are
compositions developed from those arguments.

\subsection{Divisibility-weighted and quantitative criteria}
For a finite nonempty prime set $P$, let
$h(a)=\prod_{p\in P}p^{v_p(a)}$.  The weighted condition
\[
  \sum_{a\in A}\frac{h(a)}{a(2^{h(a)}-1)}<\infty
  \tag{W}\label{eq:weighted-return}
\]
implies irrationality of $\sum_{a\in A}(b^a-1)^{-1}$ for every infinite
$A$ and every integer $b\ge2$.  The conclusion is hereditary under passage
to infinite subsets.  A fixed-base version replaces $2$ by $b$ in the
condition; arbitrary nested divisibility chains have the same criterion
with $h(a)$ the largest chain element dividing $a$.

Here the decisive estimate is a finite orbit average.  If $g=(a,Q)$, then
\[
 \frac1T\sum_{m=1}^T
 \frac{b^{Qm\bmod a}-1}{b^a-1}
 \le \frac{g}{a(b^g-1)}+\frac1{T(b^g-1)}.
\]
The second term cannot be discarded before summing over $a$.  Choose $Q$
to freeze the finite prefix and all small $P$-parts.  Average again over
$T=2^j$, $M\le j<2M$.  The incomplete-period contribution from small
$P$-parts is at most $2QW/M$, where $W$ is the sum in
\eqref{eq:weighted-return}; the large-part contribution is suppressed by
the exponential denominator.  This produces arbitrarily small positive
shifted-atom displacements.  The integral lattice gap then gives
irrationality, and the binary atom comparison gives every larger base.
The supplement gives all truncation estimates and an explicit
reciprocal-divergent host with gaps $O(\log\log a)$.

There is also a quantitative necessary condition for rationality.  Put
$H_A(x)=\sum_{a\in A,\,a\le x}1/a$, let $T_0=2$,
$T_{j+1}=2^{T_j}$, and set $\ell(x)=\min\{j:x\le T_j\}$.
For a reduced rational value $p/q$ at base $b$, write $q=q_s q_c$, where
every prime factor of $q_s$ divides $b$ and $(q_c,b)=1$.  Choose
$0\le c<q_s$ with $cq_c\equiv p\pmod {q_s}$, and put
$\lambda=1-c/q_s$.  Then
\[
 \liminf_{x\to\infty}\frac{H_A(x)}{\ell(x)}\ge\frac\lambda2.
 \tag{Q}\label{eq:logstar-mass}
\]
Indeed, at suitable multiples of the order of $b$ modulo $q_c$, rationality
forces displacement at least $\lambda$.  Applying the finite orbit estimate
after freezing the prefix through $M$, and averaging $J$ dyadic lengths,
gives, with $X=L2^{J-1}$,
\[
 J\lambda < (J+2L)\bigl(H_A(X)-H_A(M)\bigr)+4.
\]
Taking $J$ a sufficiently large multiple of $L$ forces a mass increment
arbitrarily close to $\lambda$.  The LCM bound allows this at successive
scales separated by two tower steps, giving \eqref{eq:logstar-mass}.
Consequently $H_A(x)=o(\ell(x))$ implies irrationality at every integer
base, even when the reciprocal sum diverges.  The constant is a rational
phase gap; no positive ordinary or logarithmic density is asserted.

\subsection{Positive fractional divisor covers}
Let $F_j$ be finite sets of positive integers and write
$f_{F_j}(n)=\#\{a\in F_j:a\mid n\}$.  Suppose that
$0<\alpha_j\le1$ and $c_{j,d}\ge0$ satisfy
\[
 f_{F_j}(n)^{\alpha_j}\le\sum_{d\mid n}c_{j,d},\qquad
 C_j=\sum_{d\ge1}\frac{c_{j,d}}d.
\]
Set $B_j=2^{\alpha_j}$.

\begin{theorem}[Variable-exponent positive covers]
\label{thm:variable-fractional-cover}
If
\[
 \sum_{j\ge1}C_j2^{j\alpha_j}
           \frac{B_j}{(B_j-1)^2}<\infty,
\]
then every infinite subset of $\bigcup_jF_j$ has irrational Mersenne
subseries at every integer base.  The exponents may tend to zero.
\end{theorem}
\begin{proof}
Let $U_F(n)=\sum_{r\ge1}2^{-r}f_F(n+r)$ and
$V_j(n)=\sum_{r\ge1}B_j^{-r}\sum_{d\mid n+r}c_{j,d}$.
Subadditivity gives $U_{F_j}(n)^{\alpha_j}\le V_j(n)$.
For each positive integer $L$, exact divisibility frequencies give
\[
 \lim_{X\to\infty}\frac1X\sum_{m=1}^X V_j(Lm)
 =\sum_d\frac{c_{j,d}(d,L)}{d(B_j^{(d,L)}-1)}
 \le\frac{C_j}{B_j-1}.
\]
Fix $\varepsilon>0$ and put $t_j=\varepsilon2^{-j}$.
The finite-average divisibility bound $(L+r)/d$ gives the summable
dominator
\[
 \sum_{j>J}C_jt_j^{-\alpha_j}
 \left(\frac L{B_j-1}+\frac{B_j}{(B_j-1)^2}\right).
\]
It justifies passing the average through all three sums.  The limiting
average of $\sum_{j>J}t_j^{-\alpha_j}V_j(Lm)$ is therefore at most
$\sum_{j>J}C_jt_j^{-\alpha_j}/(B_j-1)$, which tends to zero as
$J\to\infty$.  Choose $J$ making it less than one, and then choose $L$
divisible by every member of $\bigcup_{j\le J}F_j$.
Arbitrarily large $m$ make the displayed sum less than one.  At such
points every remaining $U_{F_j}(Lm)<t_j$, so their total is less than
$\varepsilon$.  The frozen atoms have zero displacement.  Infinitude
gives strictly positive displacement, and the integral lattice argument
excludes rationality.  The binary comparison transfers the returns to
all integer bases.
\end{proof}

The fixed-exponent criterion $\sum_j C_j<\infty$ follows by regrouping
finite frames into blocks with exponentially decreasing costs.  For a
complete divisor frame $F=g\{d:d\mid M\}$, positive coefficients are
provided multiplicatively by
$\beta_\alpha(p^e)=(e+1)^\alpha-e^\alpha$.  When $M$ is squarefree the
optimal cost is
\[
 \frac1g\prod_{p\mid M}\left(1+\frac{2^\alpha-1}{p}\right).
\]
This yields moving frames with divergent reciprocal mass.  Choosing
$\alpha_j=g_j^{-1/3}$ and prime blocks whose reciprocal product is about
$\exp(g_j^{1/3})$ makes the variable criterion summable when
$g_j\ge16^j$; every fixed exponent fails for this particular frame
decomposition.  Cover-independent periodic means, recorded with
\texttt{ErdosProblems/Erdos257/VariableExponentCoverSeparation.md},
exclude every alternative fixed-$\alpha$ positive cover and every
finite-prime weighted mass on an explicit support $A^\dagger$ that still
satisfies Theorem~\ref{thm:variable-fractional-cover}; a pairwise-coprime
reciprocal-divergent family defeats every positive cover.  A later ordinary
enlargement, recorded with
\texttt{ErdosProblems/Erdos257/StrengthenedVariableExponentCover.md},
weakens the displayed cost to one inverse power of $2^{\alpha_j}-1$ and
exhibits a squarefree support $A^\star$ outside every old variable-exponent
cover and every finite-prime weighted class, while every infinite subset still
has irrational subseries at every base.  Those arguments enlarge the
variable-exponent remainder.  They do not replace
Theorem~\ref{res:reciprocal-support}, which remains the Lean-checked
flagship and is the coprimality-free statement Erd\H{o}s already made.
Positivity matters: the divisor expansion for
$F=\{2,3\}$ has coefficient $2^\alpha-2<0$ at $6$ when $\alpha<1$.

\paragraph{The prime-power regime.}
A separate ordinary argument proves irrationality for every infinite
subset of the prime powers, at every integer base, including fixed
dilations and finite modifications.  Its analytic input is
Tao--Ter\"av\"ainen~\cite[Theorem~3.1]{taoteravainen2025}.
The proof uses the reciprocal mass of the selected primes as its scale,
so arbitrarily slow divergence is retained; higher prime powers produce
errors only on prime-square events.  The equidistribution, small-prime,
progression and exceptional-set hypotheses are checked in the supplement.
The full-prime theorem and the authors' stated full-prime-power extension
are theirs; no priority claim is made for the returned thinning argument.

\subsection{A finite deadline for the actual repair problem}
For the actual greedy support of $4/9$, let
$Q_N=\lfloor4\cdot2^N/9\rfloor-P_N$, with $P_N$ its integral Lambert
prefix.  Under membership, the full tail identity and divisor-pair bound
give $0\le Q_N\le2\sqrt N+4$ at every rank.  Put
$s=\lfloor\sqrt K\rfloor$ and $T=2s+12$.  Strict increase at every step
from $K$ through $K+T-1$ would imply $Q_{K+T}\ge T$.  But
$K+T<(s+4)^2$ gives $Q_{K+T}<T$, a contradiction.  Combining this with
the existing cofinal-repair consumer proves the exact equivalence
\[
 \frac49\in\Ach
 \quad\Longleftrightarrow\quad
 \forall K\ \exists N\in[K,K+2\lfloor\sqrt K\rfloor+12):
 Q_{N+1}\le Q_N.
\]
One strictly increasing window would be a finite nonmembership
certificate.  Exact finite replay through rank $4096$ finds no such
certificate and proves no all-depth statement.
The same square-root window, and cofinal nonincreases of $Q_N$, already
characterise membership for every nonnegative real target, with no rationality
or periodic-digit hypothesis; see
\texttt{ErdosProblems/Erdos257/GreedyRepairCriterion.lean} and the packet row
\texttt{generic\_real\_target\_square\_root\_repair\_criterion}.  Restricting the
display to $4/9$ hid that scope.  Every positive power $\varepsilon$ can
replace the square-root envelope, recorded with
\texttt{ErdosProblems/Erdos257/SubpowerRepairWindows.md}; that refinement does
not prove $1/2$ or $1/21$ membership.

Fixed multiplicative schedules have a different fate.  Eighty exactly
certified selected anchors through rank $727$, together with Dirichlet's
theorem, give infinitely many prime-cofactor increases of at least seven
at multiplier $120$ and at least two at multiplier $420$.  In particular,
the proposed modulus-$420$ tetraprime repair inequality is false despite
its earlier finite-window evidence.  This leaves unrestricted cofinal
repairs open.  At $1/21$, combining the actually selected exponents $5$
and $7$ gives $60$ favorable phase classes modulo $105$ and strengthens
the conditional upper-density bound for divergent quotient rows from
$1/5$ to $2/7$.  Using instead the twelve selected exponents
$5,7,8,9,10,18,20,24,28,42,45,60$ gives exactly $1011$ favorable classes
modulo $1260$.  The corresponding larger restricted divergence set has
conditional upper density at least $337/840>2/5$.  Its complete phase mask
and an independent rational replay of the selected prefix are supplied with
the synthesis.  The older $2/7$ bound concerns the two-anchor restricted set.
It proves neither target's membership.  The complete
phase certificate and the exact recurrence identities are in the
\href{https://github.com/wcook04/plectis-lean-erdos249-257/blob/main/docs/research/erdos257-eight-return-synthesis.md}{synthesis supplement}.


\section{Finite-support denominator periods}\label{sec:period}

A finite support has a rational value, so for a finite support the question is
not irrationality but arithmetic: how large the denominator is, and how the
base sits inside it.  Theorem~\ref{res:period} answers the second exactly and,
under its stated hypothesis $\lcm(F)\ge2$, gives a strict lower bound for the
first.  The omitted boundary is genuine: at $b=2$ and $F=\{1\}$ the two
quantities are both~$1$.

Six instances at $b=2$, which also show what the hypothesis
$\lcm(F)\ge2$ is for:
\begin{center}
\begin{tabular}{@{}lllll@{}}
\toprule
$F$ & $x_F(2)$ & $D_F$ & $\lcm(F)$ & $\operatorname{ord}_{D_F}(2)$\\
\midrule
$\{1\}$   & $1$      & $1$   & $1$  & $1$\\
$\{2\}$   & $1/3$    & $3$   & $2$  & $2$\\
$\{2,3\}$ & $10/21$  & $21$  & $6$  & $6$\\
$\{1,2,3\}$ & $31/21$  & $21$  & $6$  & $6$\\
$\{2,6\}$ & $22/63$  & $63$  & $6$  & $6$\\
$\{4,6\}$ & $26/315$ & $315$ & $12$ & $12$\\
\bottomrule
\end{tabular}
\end{center}
The last two columns agree in every row, which is the order statement.  In the
worked case $F=\{1,2,3\}$, the rational sum is $1+1/3+1/7=31/21$ and
$2^6\equiv1\pmod {21}$, whereas no smaller positive exponent gives~$1$.
The
first row is the only finite nonempty support with $\lcm(F)=1$, and the only
one on which $\lcm(F)<D_F$ fails; that is what the hypothesis $\lcm(F)\ge2$
excludes.

The content is a noncollapse statement.  Clearing denominators over
$b^{\lcm(F)}-1$ makes the period at most $\lcm(F)$ immediately; what is not
immediate is that cancellation in the numerator cannot bring it below.  A
divisibility-maximal exponent $n\ge2$ contributes a cyclotomic prime power
$\ell^e$ with $e=v_\ell(b^n-1)$ on which $b$ has order exactly $n$.  That
valuation is uniquely minimal in the finite sum, so it survives reduction.
The unsigned theorem is this argument with every sign $+1$.  The signed
unit-coefficient form, and the identification $\operatorname{ord}_{\ell^e}(b)=n$
when $\ell\mid\Phi_n(b)$, are Lean-checked in
\texttt{ErdosProblems/Erdos257/SignedFinitePeriodNoncollapse.lean}; the ordinary
write-up is
\texttt{ErdosProblems/Erdos257/FiniteDenominatorRealisation.md}.  Prime powers rather
than primes is not a technicality.  At $(b,n)=(2,6)$ no prime divisor of
$2^6-1=63$ has order $6$, one of the exceptional cases in Zsigmondy's
theorem, while $9$ does, and it
is the prime power that carries the local witness for $F=\{2,6\}$.  Order $6$
has two distinct mechanisms: $\operatorname{ord}_9(2)=6$ for the maximal
exponent $6$ in $F=\{2,6\}$, while $\operatorname{ord}_{21}(2)=\operatorname{lcm}(2,3)=6$
for $F=\{2,3\}$.  The growth clause
then follows formally, since the order divides $\ph(D_F)<D_F$, where $\ph$ is
Euler's totient.

Two boundaries.  This is an unconditional statement about every finite
support, not a bounded table of examples, and not an implication with an open
hypothesis; but it settles no infinite support, and no limit of it does.
Denominator-period control is part of the classical method behind
Erd\H{o}s's 1948 argument~\cite{erdos1948}.  Whether this exact sharp form
appears in the literature has not been assessed, and no novelty is claimed.

\section{Rational values and scaled tails}\label{sec:forced}

Section~\ref{sec:map} below lists supports for which irrationality is known.
This section is the other half, and the half that meets
Problem~\ref{res:problem} rather than sampling it: statements that hold for
\emph{every} infinite support whose value is rational.  The base is $2$
throughout, and each statement carries its own hypotheses.

\paragraph{Integral scaled tails, and unboundedness.}
Binary long division turns a hypothetical rational value into an integer
recurrence.  Fix an integer $v\ge1$ and a sequence $f\colon\Npos\to\N$ of
coefficients.  Call an integer sequence $u\colon\N\to\Z$ a
\emph{tempered scaled-tail sequence} for $f$ with multiplier $v$ if
\[
  u(n+1)=2u(n)-v\,f(n+1)\quad\text{for every }n\ge0,
  \qquad\text{and}\qquad u(n)=o(2^{n}).
\]

The recurrence is one step of long division in base~$2$: double the remainder,
then pay out the next coefficient.  The growth condition is what pins the
solution down, since two solutions of the recurrence differ by $C\,2^{n}$ for a
constant $C$, and only $C=0$ survives $u(n)=o(2^{n})$.  For any coefficient
sequence with $f(n)\le n$, and $\scA_A(n)\le d(n)\le n$ qualifies, $d$ being
the divisor function, the series $\sum f(n)2^{-n}$ is rational exactly when
a tempered scaled-tail sequence exists for some $v\ge1$, and every such
sequence is then
the scaled tail
$u(n)=v\sum_{j\ge1}f(n+j)2^{-j}$, so the orbit is unique rather than merely
available
(\mword{Erdos249257/GenericTailOrbitRigidity.lean}{426}{binaryCoeffSeries_rational_iff_exists_temperedBinaryOrbit}{checked}).

The next theorem exhibits such an orbit for $X_A$ itself, with the coefficient
sequence shifted past the power of $2$ in the denominator.

\begin{theorem}[unbounded scaled-tail states]\label{res:unbounded}
Let $A\subseteq\Npos$ be infinite with $X_A=p/(2^cv)$, $v\ge1$.  Then there is
$u\colon\N\to\Npos$ with
\[
  u(n)=v\sum_{j\ge1}\scA_A(c+n+j)2^{-j}
\]
satisfying the exact recurrence
\[
  u(n+1)+v\,\scA_A(c+n+1)=2u(n),
  \qquad u(n)\equiv p\,2^{n}\ (\mathrm{mod}\ v),
\]
and $u$ is unbounded.
\end{theorem}

\noindent Checked as
\mword{Erdos249257/RationalSupportCarrySkeleton.lean}{2383}{exists_unbounded_shifted_odd_tail_nat_state_of_support_fraction}{Lean}.
The arithmetic clauses are the integral lattice, the residue class, and the
exact recurrence.  Unboundedness itself holds for every infinite support, by
taking common multiples of arbitrarily many selected exponents; see
\texttt{ErdosProblems/Erdos257/ElementarySupportConstraints.md}.  It is not a
rationality restriction.

\paragraph{Divisor coverage cannot have long gaps.}
Say that $\scA_A$ has a \emph{zero window} of length $h$ at $N$ if
$\scA_A(N+1)=\dots=\scA_A(N+h)=0$, that is, if no element of $A$ divides any
of $h$ consecutive integers.  At $A=\{2,3\}$, for instance, $\scA_A$ vanishes
exactly on the integers prime to $6$, so every zero window has length at
most $1$: one of any two consecutive integers is even.

\begin{theorem}[sublogarithmic zero windows]\label{res:sublog}
Let $A\subseteq\Npos$ be nonempty with $X_A=p/(2^cv)$, $v\ge1$.  For every
$\varepsilon>0$ there is a constant $B$, depending on
$A,p,\varepsilon,c$ and $v$, such that every zero window of $\scA_A$ at
$c+N$ has length
\[
  h\le\varepsilon\log_2(N+1)+B .
\]
For fixed $A,p,c,v$ and $\varepsilon$, the same $B$ works uniformly in $N$
and $h$; no common constant over supports or numerators is asserted.
\end{theorem}

\noindent Checked as
\mword{Erdos249257/SublogDivisorCoverage.lean}{392}{supportCoeffZeroWindow_length_le_eps_logb_add}{Lean}.
The Lean bound is true.  Independently of rationality, every nonempty support
already has zero windows of length at most $\min A-1$; see
\texttt{ErdosProblems/Erdos257/ElementarySupportConstraints.md}.  The
sublogarithmic statement is therefore not a substantial extra restriction on
rational counterexamples.

\paragraph{A reciprocal-sum lower bound.}
Let $\operatorname{ord}_v(2)$ denote the multiplicative order of $2$ modulo an
odd $v>1$, and write the reciprocal mass of $A$ for $\sum_{a\in A}1/a$.

\begin{theorem}[reciprocal mass bound]\label{res:mass}
Let $A\subseteq\Npos$ be such that $\sum_{a\in A}1/a$ converges, and suppose
$X_A=p/(2^cv)$ with $v>1$ odd and $\gcd(p,v)=1$.  Then
\[
  \sum_{a\in A}\frac1a\;\ge\;\frac1{\operatorname{ord}_v(2)} .
\]
\end{theorem}

\noindent Checked as
\mword{Erdos249257/RationalSupportCarrySkeleton.lean}{1480}{one_div_oddOrder_le_reciprocalMass_of_support_fraction}{Lean}.
A finite support already supplies an instance: $A=\{2,3\}$ has $X_A=10/21$, so
$c=0$, $v=21$, $\gcd(p,v)=1$ and $\operatorname{ord}_{21}(2)=6$, and the bound
reads $\tfrac12+\tfrac13\ge\tfrac16$.
The summability hypothesis is doing real work: a rational value with fixed odd
denominator part $v$ and a convergent reciprocal sum requires mass at least
$1/\operatorname{ord}_v(2)$.  The statement does not apply when the reciprocal
sum diverges, and it gives no positive lower bound uniform in $v$; it
constrains the pair (support, denominator) jointly.

The critical dyadic case has a different endpoint.

\begin{theorem}[dyadic reciprocal-mass endpoint]\label{res:dyadicmass}
Let $A\subseteq\Npos$ be infinite and suppose $X_A=p/2^c$ for some
$p\in\Z$ and $c\in\N$.  Then the reciprocal support terms are not summable,
or
\[
  \sum_{a\in A}\frac1a>1.
\]
\end{theorem}

\noindent Checked as
\mword{Erdos249257/RationalSupportCarrySkeleton.lean}{2210}{dyadic_support_fraction_reciprocalMass_diverges_or_gt_one}{Lean}.
The two alternatives are an exact necessary consequence of a dyadic rational
value.  They do not contradict rationality by themselves: a contradiction
requires separate proofs that the reciprocal terms are summable and that the
mass is at most $1$.  The formal nonsummability alternative also does not, by
itself, assert a particular asymptotic law for the partial sums.

\paragraph{A Boolean--M\"obius characterisation.}
The preceding theorems give necessary conditions.  The next statement is an
equivalence, and is the sharpest description of rational-valued supports the
development has.  It removes the support from the description altogether, which
is possible because the support can always be read back off the incidence
sequence: with
$\mu$ the M\"obius function, $*$ Dirichlet convolution
$(g*h)(n)=\sum_{d\mid n}g(d)h(n/d)$ and $\mathbf 1_A$ the indicator of $A$, the
identity $\scA_A=\mathbf 1_A*1$ inverts to $\mu*\scA_A=\mathbf 1_A$.

For $p\in\Z$, $q\ge1$ and $U\colon\N\to\Z$, set
\[
  Q_U(0)=0,\qquad
  Q_U(n)=\frac{2U(n-1)-U(n)}q\quad(n\ge1).
\]
Call $U$ an \emph{admissible Boolean--M\"obius scaled-tail sequence} for
$(p,q)$ if
\[
\begin{gathered}
 U(0)=p,\quad U(N)>0,\quad
 U(N)\le q\bigl(2\sqrt N+4\bigr)\quad(N\ge0),\\
 q\mid 2U(N)-U(N+1)\quad(N\ge0),\qquad
 (\mu*Q_U)(n)\in\{0,1\}\quad(n\ge1).
\end{gathered}
\tag{2}\label{eq:bmc}
\]
The divisibility condition makes $Q_U$ integral; the last condition says that
M\"obius inversion of the quotient is the indicator of a set.

\begin{theorem}[Boolean--M\"obius scaled-tail correspondence]\label{res:bmc}
Let $p\in\Z$ and $q\ge1$.  There exists a support $A$ with $0\notin A$, with
some positive element, and with $X_A=p/q$, if and only if there exists an
admissible Boolean--M\"obius scaled-tail sequence $U$ for $(p,q)$.  In that
case the support is recovered as
\[
  A=\{n:(\mu*Q_U)(n)=1\},
\]
with $Q_U$ as in~\eqref{eq:bmc}.
\end{theorem}

\noindent Checked as
\mword{Erdos249257/BooleanMobiusCarry.lean}{949}{exists_normalized_support_fraction_iff_exists_booleanMobiusCarry}{Lean}.
The finite example $A=\{2,3\}$ makes the definition concrete.  Here
$p/q=10/21$, and
\[
 Q_U(1),\ldots,Q_U(6)=0,1,1,1,0,2,\qquad
 U(0),\ldots,U(6)=10,20,19,17,13,26,10.
\]
M\"obius inversion gives
$(\mu*Q_U)(1),\ldots,(\mu*Q_U)(6)=0,1,1,0,0,0$, the indicator of
$\{2,3\}$ through that range.  This finite calculation illustrates the
correspondence; the admissible scaled-tail sequence itself is an infinite
object, and the finite support is not a counterexample to
Problem~\ref{res:problem}.

One boundary on the equivalence.  The support it produces is not required to be
infinite, and a finite support supplies an admissible scaled-tail sequence for
its own value, so existence of such a sequence is not by itself a
counterexample to Problem~\ref{res:problem}.
What the equivalence changes is the search space, not the difficulty.

\paragraph{Combined constraints on a rational counterexample.}
Taken together: a counterexample to Problem~\ref{res:problem} would have an
unbounded scaled-tail sequence obeying an exact linear recurrence and
sublogarithmic divisor-coverage gaps.  If its reciprocal sum converged and its
reduced odd denominator part were greater than~$1$, it would also satisfy
Theorem~\ref{res:mass}; at a dyadic rational value it would instead satisfy
Theorem~\ref{res:dyadicmass}.  Theorem~\ref{res:bmc} reconstructs every
rational value, but by itself does not force the reconstructed support to be
infinite.  These qualifications matter: the statements have different
hypotheses, and no jointly contradictory combination has been proved.

\section{Representative known irrational supports}\label{sec:map}

The following table is representative, not exhaustive.  It groups the
displayed supports by six mechanisms rather than suggesting that its rows
classify all known cases.

\begin{table}[htbp]
\centering
\small
\begin{tabular}{@{}
  >{\raggedright\arraybackslash}p{0.29\linewidth}
  >{\raggedright\arraybackslash}p{0.15\linewidth}
  >{\raggedright\arraybackslash}p{0.48\linewidth}@{}}
\toprule
Support $A$ & Bases & Mechanism and authority \\
\midrule
\multicolumn{3}{@{}l}{\emph{LCM-prefix orbit means and an integral tail gap}}\\
arbitrary infinite $A$ with $\sum_{a\in A}a^{-1}<\infty$ & every $b\ge2$ &
  Theorem~\ref{res:reciprocal-support}; Lean-checked all-base theorem obtained
  by transferring the independently reconstructed binary close returns, with
  exact boundary in Section~\ref{sec:reciprocal-support} \\
\addlinespace
\multicolumn{3}{@{}l}{\emph{Classical full support, and base change into it}}\\
all of $\Npos$ & every $b\ge2$ & Erd\H{o}s 1948~\cite{erdos1948};
  \mword{Erdos249257/CertificateKernel.lean}{9045}{irrational_erdosSupportSeries_univ}{formalised here} \\
multiples of a fixed $d$ & every $b\ge2$ & dilation: the multiples series at
  base $b$ \emph{is} the full-support series at base $b^d$
  (\mword{Erdos249257/CertificateKernel.lean}{9103}{irrational_erdosSupportSeries_multiples}{checked}) \\
\addlinespace
\multicolumn{3}{@{}l}{\emph{Periodicity of the coefficient sequence}}\\
eventually periodic & every $b\ge2$ &
  \mword{Erdos249257/CertificateKernel.lean}{12811}{irrational_ratWeightSeries_eventuallyPeriodic}{checked here};
  Luca--Tachiya prove the nonnegative purely-periodic case
  \cite[Theorem~1, p.~139; proof pp.~149--150]{lucatachiya2017};
  finite rational prefixes give the infinite eventual case \\
a residue class; the odd numbers & every $b\ge2$ & special cases of the row
  above (\mword{Erdos249257/CertificateKernel.lean}{11672}{irrational_erdosSupportSeries_residueClass}{residue class},
  \mword{Erdos249257/CertificateKernel.lean}{11686}{irrational_erdosSupportSeries_odd}{odd});
  Luca--Tachiya's Example~2 strengthens the odd row to joint linear
  independence of every finite divisor-convolution ladder, also for negative
  integer bases with absolute value greater than one
  \cite[Example~2, p.~140]{lucatachiya2017} \\
\addlinespace
\multicolumn{3}{@{}l}{\emph{Denominator gaps along a sparse support}}\\
factorials $\{n!\}$ & every $b\ge2$ &
  \mword{Erdos249257/CertificateKernel.lean}{6035}{irrational_erdosSum_factorial_support}{checked here} \\
powers of two $\{2^k\}$ & every $b\ge2$ &
  \mword{Erdos249257/CertificateKernel.lean}{6059}{irrational_erdosSum_two_pow_support}{checked here} \\
pairwise coprime, $\sum_{a\in A}a^{-1}<\infty$ & every $b\ge2$ &
  Erd\H{o}s~\cite{erdos1968}, theorem on p.~222;
  \mword{Erdos249257/CertificateKernel.lean}{10776}{irrational_erdosSupportSeries_pairwise_coprime}{formalised here} \\
\addlinespace
\multicolumn{3}{@{}l}{\emph{Multiplicative closure: Chowla--Erd\H{o}s, refined}}\\
$F_s(E)$ and the monomial images
  $\{j n^i:n\in F_s(E)\}$ & $b=q^j$ with $|q|^L\le s$,
  $L=\lcm(1,\dots,\ell)$ &
  Duverney--Tachiya~\cite[Cor.~1.2, author-preprint p.~4]{duverneytachiya};
  linear independence, not only irrationality; not formalised here \\
$s$-free positive integers, with or without $1$ &
  $b=q^j$, $2\le |q|\le s$, $j\ge1$ &
  \emph{ibid.}, with $E$ the primes and $\ell=i=1$;
  deleting $1$ changes the value by $1/(b-1)\in\Q$ \\
squarefree & every $b=2^j$, $j\ge1$ &
  Duverney--Tachiya~\cite[Cor.~1.2 and Ex.~1.1, author-preprint p.~4]{duverneytachiya};
  joint linear independence for every finite set of such bases; not formalised here \\
coprime to a fixed $N$; sums of two squares & every $b\ge2$ &
  \emph{ibid.}, Examples 1.3 and 1.2; not formalised here \\
\addlinespace
\multicolumn{3}{@{}l}{\emph{Analytic}}\\
primes & $b=2$ proved & Tao--Ter\"av\"ainen~\cite{taoteravainen2025},
  Thm.~1.3, p.~4; proof pp.~44--56; not formalised here \\
primes, $b\ge3$; prime powers & --- & \emph{ibid.}, asserted as modifications
  with details omitted in the cited version \\
\midrule
arbitrary infinite $A$ & --- & \textbf{Open} (Problem~\ref{res:problem}) \\
\bottomrule
\end{tabular}
\end{table}

\noindent In the sentence immediately following the p.~222 theorem,
Erd\H{o}s says that pairwise coprimality can be removed by a more complicated
argument, but he does not give that argument; p.~226 repeats that boundary.
His printed theorem already quantifies over every integer base $t\ge2$, so the
base generality of the first row is his.  The first row records a complete
Lean-checked proof of the conclusion he stated and withheld, namely the same
theorem with pairwise coprimality dropped.  The sparse-support row separately
records his printed pairwise-coprime theorem and its formalisation.  Removing
coprimality is the contribution, and it carries no priority claim.

Under the additional hypothesis that the pairwise-coprime exponents are odd,
Duverney, Suzuki and Tachiya strengthen this conclusion: Corollary~1 gives
linear independence of (1) and three associated plus/minus reciprocal sums
at every integer base of absolute value greater than one
\cite[Corollary~1]{duverneysuzukitachiya2019}.  This is a stronger conclusion
on a narrower support class; it does not remove the hypothesis from the row
above.

\noindent Five remarks on the table: one on containment between the
mechanisms, three on the literature rows, and one on what ``checked here''
means.

For the full-support row, the weighted certificate theorem is
\mword{Erdos249257/CertificateKernel.lean}{8328}{irrational_erdosSum_full_support}{the exact checked statement};
the support-series formulation displayed in the table is its later corollary.

\emph{The mechanisms are not nested, and one non-containment is proved.}  The
base-$b$ full-support theorem is exactly the $A=\Npos$ statement; the
multiples rows genuinely specialise it after a base change, but the sparse
rows do not, and this is not merely an artefact of how they were proved.  The
denominator-gap criterion behind the factorial and power-of-two rows
\emph{provably cannot} reach the full support: the prefix lcm grows too slowly
for its hypothesis to hold
(\mword{Erdos249257/CertificateKernel.lean}{6272}{lcm_gap_hypothesis_fails_full_support}{checked}).
Conversely the analytic method of~\cite{taoteravainen2025} reaches the primes,
which are neither eventually periodic nor pairwise coprime with summable
reciprocals.  No list contains another, and their union does not exhaust the
infinite supports.

\emph{One theorem of Duverney and Tachiya generates several rows.}
Their refinement of the Chowla--Erd\H{o}s method
~\cite[Cor.~1.2, author-preprint p.~4; proof pp.~9--11]{duverneytachiya}
proves, for a pairwise coprime sequence $E$ of polynomial growth and the set
$F_s(E)$ of products of its members with exponents below $s$, that
$1$ and the values $\sum_{n\in F_s}(q^{jn^i}-1)^{-1}$ are linearly independent
over $\Q$ whenever $|q|^{L}\le s$, with $L=\lcm(1,\dots,\ell)$ and $\ell$
bounding the exponent $i$.  This is a
row-generating theorem, not an isolated example.

With $E$ the primes,
$s=\infty$ returns the full support at every base, while $s=2$ returns the
squarefree support.  For \(s=2\), the constraint forces \(\ell=1\) and
\(|q|\le2\), but the free exponent \(j\) gives every base \(b=2^j\).
Thus this theorem excludes bases that are not powers of \(2\), rather than all
bases \(b\ge3\); it also excludes higher monomial degrees \(i\ge2\).
Their conclusion is stronger than irrationality, being linear independence of
the whole finite family across the chosen exponents \(j\).

\emph{The classical full-support value now has a digit-level refinement.}
At base \(2\), Campbell writes the Erd\H{o}s--Borwein constant in the
equivalent forms
\[
  E=\sum_{n\ge1}\frac1{2^n-1}
   =\sum_{n\ge1}\frac{d(n)}{2^n}
\]
and proves that the binary block \(11\) occurs infinitely often in its
base-\(2\) expansion
\cite[Theorem~1, p.~12; proof pp.~12--24]{campbell2026}.
This adds genuine digit-distribution information to the full-support row, but
it neither proves normality nor addresses an arbitrary infinite support
\(A\), so it does not change the open status of Problem~\ref{res:problem}.

\emph{The two analytic rows do not have the same standing.}  Theorem 1.3
on p.~4, proved in Section~5 on pp.~44--56,
of~\cite{taoteravainen2025} proves the prime-support case at base~$2$, the
series there being $\sum_{n\ge1}\omega(n)/2^n$.  The extension to every
integer base, and the prime-power support, which those authors themselves
identify with Problem~\ref{res:problem}, are asserted by remark, with the
modifications explicitly left to the reader.  The table keeps the three apart,
and only the first is proved.

\emph{Formalisation is not priority.}  Rows marked ``checked here'' are Lean
statements accepted by the pinned kernel.  For the full support that is a
formalisation of Erd\H{o}s; Luca and Tachiya's RIMS paper already proves the
nonnegative purely-periodic case; a finite rational-prefix correction gives
the eventual-periodic extension used here
\cite[Theorem~1, p.~139; proof pp.~149--150]{lucatachiya2017}, while
Theorem~A restates the broader signed purely-periodic theorem without
reproducing its earlier proof.  No priority is claimed anywhere in this table.

\begin{remark}[signed periodic weights: a gap in the method, not in the
mathematics]\label{res:signed}
Dropping nonnegativity weakens the conclusion available, and it is worth being
exact about where.  For $u\colon\Npos\to\Z$ periodic and
$b\ge2$, what is checked here is a \emph{dichotomy}: the series
$\sum_{n\ge1}u(n)/(b^n-1)$ is irrational, \emph{or} some power of $b$ times
its value is an integer
(\mword{Erdos249257/CertificateKernel.lean}{14175}{irrational_or_bpow_mul_eq_intCast_intWeightedErdosSeries_periodic}{checked}).
For a nonnegative support indicator the terminating alternative is excluded by
the rows above; for mixed signs this argument does not exclude it.

The terminating alternative never in fact occurs.  Theorem~A in Luca and
Tachiya's 2017 RIMS paper explicitly restates their earlier Theorem~1.1 in the
exact form needed here: $\sum_{n\ge1}a_n/(q^n-1)$ is irrational for every
integer $q$ with $|q|>1$ and every nonzero purely periodic integer sequence
$(a_n)$ \cite[Theorem~A, p.~139]{lucatachiya2017}.  Thus the disjunction above
is a consequence of the chosen argument, not a feature of the problem: the second
branch is empty, and their author restatement is what shows it.  The RIMS paper
does not reproduce the earlier proof or verify the broader eventual-rational
formulation, so neither is claimed here.  The local statement retains only
the interest of being an independently checked argument with a visible
boundary.  Section~\ref{sec:squarefree} records a second, sharper instance of
the same phenomenon.
\end{remark}

\section{Squarefree support and a coordinate-dependent obstruction}
\label{sec:squarefree}

Let $\Asf=\{d\ge2: d\text{ squarefree}\}$: a support of density $6/\pi^2$, far
denser than the primes, and not periodic.  Its values at all power-of-two bases
are jointly linearly independent with \(1\), by a theorem of Duverney and
Tachiya recalled below.  This section establishes that two block-certificate
arguments used here cannot reach that support at any even base, and that the
reason lies in a normalisation rather than in the value.

\begin{theorem}[squarefree divisor incidence]\label{res:sfcount}
For every $n\ge1$ the number of squarefree divisors of $n$ is $2^{\omega(n)}$,
where $\omega(n)$ is the number of distinct prime factors, and hence
\[
  \scA_{\Asf}(n)=2^{\omega(n)}-1 .
\]
In particular $\scA_{\Asf}(n)$ is \emph{odd} for every $n\ge2$.
\end{theorem}

\noindent The count is
\lword{Erdos257/SquarefreeSupportIncidence.lean}{94}{card_squarefreeDivisors}{checked},
the incidence formula is
\lword{Erdos257/SquarefreeSupportIncidence.lean}{111}{squarefreeIncidence_eq}{checked},
and the parity conclusion is
\lword{Erdos257/SquarefreeSupportIncidence.lean}{140}{odd_squarefreeIncidence}{checked}.
The proof is the bijection between squarefree divisors of $n$ and subsets of
its prime factors
(\lword{Erdos257/SquarefreeSupportIncidence.lean}{71}{squarefreeDivisors_eq_image}{Lean}),
so the count is $2^{\omega(n)}$; removing $d=1$ leaves an odd number whenever
$\omega(n)\ge1$.  At $n=12=2^2\cdot3$ the squarefree divisors are $1,2,3,6$, so
$\scA_{\Asf}(12)=3$; at $n=30$ they are the eight divisors of $30$ and
$\scA_{\Asf}(30)=7$.

\subsection{The values are known at every power-of-two base}

Duverney and Tachiya's Corollary~1.2, applied with $E$ the primes, $s=2$ and
$\ell=1$, gives $L=\lcm(1)=1$ and the admissibility condition $|q|\le2$.
The exponent \(j\) in their displayed family remains arbitrary.  Thus, for
every \(h\ge1\), their Example~1.1 says that the numbers
\[
  1,\qquad
  \sum_{n\ge1}\frac{|\mu(n)|}{2^{jn}-1}\quad(1\le j\le h)
\]
are linearly independent over $\Q$
~\cite[Cor.~1.2 and Ex.~1.1, author-preprint p.~4]{duverneytachiya}.  Since
$|\mu|$ is the indicator of the squarefree integers including $n=1$, and
$\Asf$ omits only $n=1$, whose weight at base \(2^j\) is
\((2^j-1)^{-1}\in\Q\),
\[
  X_{\Asf}(2^j)
  =\sum_{n\ge1}\frac{|\mu(n)|}{2^{jn}-1}-\frac1{2^j-1},
\]
and rational translation preserves the joint linear independence with \(1\).

\begin{corollary}[joint power-of-two-base theorem]\label{res:sfsettled}
For every \(h\ge1\), the \(h+1\) numbers
\[
  1,\quad X_{\Asf}(2),\quad X_{\Asf}(4),\quad\ldots,\quad X_{\Asf}(2^h)
\]
are linearly independent over \(\Q\).  In particular, every
\(X_{\Asf}(2^j)\), \(j\ge1\), is irrational.
\end{corollary}

\noindent Two boundaries on that.  It is a citation, not a formalisation:
nothing in this development proves it.  The admissibility condition
$|q|^L\le s$ at \(s=2\) allows only \(|q|=2\), so this citation covers the
bases \(2^j\) and does not cover bases that are not powers of \(2\); this note
proves no result about those remaining values.

\subsection{Two block-certificate hypotheses have no instance}

Fix $b\ge2$ and a nonnegative coefficient sequence $f$.  The two hypotheses
used below are best printed rather than named.  For every precision $q\ge1$,
they ask for $N,K,L,C\in\N$ with $K\le L$ satisfying the common conditions
\[
\begin{gathered}
 \sum_{r=K+1}^{L}f(N+r)b^{L-r}\le C,\qquad
 \exists t\in\N:\ f(N+L+1+t)>0,\\
 q(C+N+L+2)<b^L,
\end{gathered}
\tag{3}\label{eq:block-common}
\]
together with one of the two first-block conditions
\[
\begin{array}{ll}
\text{digitwise:}& b^r\mid f(N+r)\quad(1\le r\le K),\\[2mm]
\text{weighted:}&
b^K\mid\displaystyle\sum_{r=1}^{K}f(N+r)b^{K-r}.
\end{array}
\tag{4}\label{eq:block-first}
\]
Thus the data $N,K,L,C$ may depend on $q$.  The parity of
Theorem~\ref{res:sfcount} refutes both alternatives for the squarefree
incidence sequence at every even base.

\begin{corollary}[neither divisibility-first hypothesis has an instance]\label{res:blind}
For the squarefree support $\Asf$ and every even base $b\ge2$, neither the digitwise
nor the carry-aware block-certificate hypothesis holds.  Both fail already at
precision $q=b^2$.
\end{corollary}

\begin{proof}
The hypothesis quantifies over every precision, so it is enough to exhibit one
precision at which no admissible data exist; we take $q=b^{2}$ and split on
whether the first block is empty.

Both hypotheses ask, for every precision $q$, for $N$, $K\le L$ and $C$ with a
first-block condition on $\scA_{\Asf}(N+1),\dots,\scA_{\Asf}(N+K)$, a middle
bound $\sum_{r=K+1}^{L}\scA_{\Asf}(N+r)b^{L-r}\le C$, a nonzero coefficient
beyond $N+L$, and $q(C+N+L+2)<b^{L}$.  Take $q=b^2$ and write $f=\scA_{\Asf}$.

Suppose $K\ge1$.  The digitwise condition requires $b^{r}\mid f(N+r)$ for
$1\le r\le K$.  If $N+K\ge2$, its instance at $r=K$ makes the even number
$b$ divide the odd number $f(N+K)$.  The carry-aware condition requires
$b^{K}\mid\sum_{r=1}^{K}f(N+r)b^{K-r}$; reducing modulo $b$ kills every term
but $r=K$, so the same contradiction follows when $N+K\ge2$.

The only remaining case with $K\ge1$ is $N=0$, $K=1$.  If $L=1$, then
$q(C+N+L+2)\ge3b^2>b$; if $L\ge2$, the $r=2$ term in the middle sum is
$f(2)b^{L-2}=b^{L-2}$, so $qC\ge b^L$.  Both contradict the strict size
inequality.

Suppose $K=0$, where both first-block conditions are vacuous.  If $L=0$ the
middle sum is empty and $q(C+N+2)<b^{0}=1$ is impossible.  If $L\ge1$ the
$r=1$ term gives $C\ge b^{L-1}$ whenever $N\ge1$, and again the size
inequality fails.  If $N=0$ and $L=1$, its left side is at least
$3b^2>b$; if $N=0$ and $L\ge2$, the $r=2$ term gives
$C\ge b^{L-2}$ and hence $qC\ge b^L$.  These exhaust the cases.
\end{proof}

\noindent The two formal nonexistence statements are checked directly as
\lword{Erdos257/SquarefreeSupportIncidence.lean}{274}{not_exists_carry_certificates_squarefreeSupport}{the carry-aware no-go}
and
\lword{Erdos257/SquarefreeSupportIncidence.lean}{292}{not_exists_digitwise_certificates_squarefreeSupport}{the digitwise no-go}.
The displayed proof is retained to expose the empty- and one-position-block
boundary cases rather than leaving the scope hidden behind the interfaces.

\subsection{The obstruction is a normalisation, not the value}

Corollary~\ref{res:blind} is a statement about two arguments failing on a value
that Corollary~\ref{res:sfsettled} shows to be irrational.  The question it
raises is therefore not whether the value can be reached, but what the failure
is a property \emph{of}.  It is a property of where the support starts.

Adjoin $1$ to the support and write $\Asf^{+}=\Asf\cup\{1\}$, the full
squarefree support.  Then
\[
  X_{\Asf^{+}}(b)-X_{\Asf}(b)=\frac1{b-1}\in\Q ,
\]
so the two supports pose the same irrationality question at every base, while
the divisor incidence changes from $2^{\omega(n)}-1$ to $2^{\omega(n)}$,
from odd to even at every $n\ge2$.  The parity obstruction of
Corollary~\ref{res:blind} evaporates under a shift that provably cannot change
the answer.  The shifted coefficient identity and the exact equivalence of the
two irrationality questions are checked as
\lword{Erdos257/SquarefreeSupportIncidence.lean}{314}{supportCoeff_squarefreeShiftedSupport}{$2^{\omega(n)}$ incidence}
and
\lword{Erdos257/SquarefreeSupportIncidence.lean}{335}{irrational_erdosSupportSeries_squarefreeSupport_iff_shifted}{the shift iff}.

More is true: the shifted first-block condition at base~$2$ asks for
$2^{r}\mid 2^{\omega(N+r)}$, that is $\omega(N+r)\ge r$ for $1\le r\le K$, and
a Chinese-remainder construction reserving $r$ fresh primes for each shift $r$
supplies such an $N$ for every $K$.  This is checked both as the arithmetic
block theorem
\lword{Erdos257/SquarefreeSupportIncidence.lean}{441}{exists_omega_ge_block}{for $\omega$}
and in the corresponding formal statement
\lword{Erdos257/SquarefreeSupportIncidence.lean}{485}{exists_digitwise_block_squarefreeShiftedSupport}{for the shifted incidence}.

It does \emph{not} follow that either argument certifies the shifted support:
the opening block is one of the conditions listed above, and the middle bound
and the arithmetic inequality are untouched by this observation.  What does
follow is a methodological point: an obstruction stated against a coefficient
sequence can be an artefact of the normalisation chosen for that sequence.
Before a no-go result is reported as a property of a problem, the coordinate
it is stated in should be varied by a transformation the problem is known to
be invariant under.  Here the transformation is adding one rational number,
and it removes the obstruction entirely.

The same invariance holds for every finite change, not only this one.  If
$A\mathbin{\triangle}B$ is finite, choose $M$ above all of its elements.
The two support series then have the same tail beyond $M$, while each omitted
prefix is rational.  The two directions of this argument are exactly the
checked prefix lemmas
\mword{Erdos249257/CertificateKernel.lean}{9467}{irrational_erdosSupportSeries_of_tail}{Lean}
and
\mword{Erdos249257/CertificateKernel.lean}{9476}{irrational_erdosSupportSeries_tail_of_irrational}{Lean}.
Thus $X_A(b)$ is irrational if and only if $X_B(b)$ is irrational for every
integer $b\ge2$.  The present shift is the smallest instance of a general
checked finite-change principle.

This is the second time in this note that a boundary turns out to belong to
the method rather than to the mathematics.  In Remark~\ref{res:signed} the
terminating alternative of the signed periodic dichotomy is empty, and a
theorem of Luca and Tachiya is what shows it; here a parity obstruction
survives only until the support is shifted by one element.  In each case the
correction came from outside the development, once from the literature and once
from asking what the statement was invariant under.  A formalised no-go result
carries exactly the authority of its hypotheses, and its hypotheses include the
coordinates it was written in.

\section{Achievement-set geometry and the value \texorpdfstring{$1/2$}{1/2}}
\label{sec:geometry}

Problem~\ref{res:problem} asks whether every value obtainable from infinitely
many of the weights $\wgt n=(2^n-1)^{-1}$ is irrational.  The set of all values
obtainable, from finite and infinite selections alike, is the achievement set
$\Ach=\{\sum_{n\ge1}\varepsilon_n\wgt n:\varepsilon_n\in\{0,1\}\}$, and this
section is about its geometry.

\begin{theorem}[geometry of $\Ach$]\label{res:geometry}
$\Ach$ is compact, perfect, totally disconnected and nowhere dense, and has
Lebesgue measure~$1$.
\end{theorem}

\noindent Checked as
\mword{Erdos249257/GreedyAchievementSet.lean}{656}{isCompact_mersenneAchievementSet}{compact},
\mword{Erdos249257/GreedyAchievementSet.lean}{1633}{perfect_mersenneAchievementSet}{perfect},
\mword{Erdos249257/GreedyAchievementSet.lean}{1649}{isTotallyDisconnected_mersenneAchievementSet}{totally disconnected},
\mword{Erdos249257/GreedyAchievementSet.lean}{1658}{isNowhereDense_mersenneAchievementSet}{nowhere dense},
and
\mword{Erdos249257/GreedyAchievementSet.lean}{996}{volume_mersenneAchievementSet}{of measure one}.
So $\Ach$ is a fat Cantor set: strict tail domination
$\sum_{\ell>n}\wgt\ell<\wgt n$ opens a gap at every level, while the total
measure is not lost.  The classical antecedent is Kakeya's two-page paper: he
proves perfectness for the subsum set of every absolutely convergent real or
complex series and gives the real interval criterion; under the present
strict tail inequality at every index, his result also gives nowhere density
\cite[p.~251]{kakeya1914}.  The strict inequality, the resulting distinctness
of subsums over distinct supports, and the Cantor conclusion for every integer
base are also recorded in Remark~4.1 (p.~13) of Kova\v{c} and
Tao~\cite{kovactao}; no novelty is claimed for any of these facts.  That
remark makes no metric assertion.

Strict tail domination does not determine the measure: the weights $3^{-n}$
satisfy $\sum_{\ell>n}3^{-\ell}=3^{-n}/2<3^{-n}$ and produce a null achievement
set, the base-$3$ digits-in-$\{0,1\}$ Cantor set.  So the measure-one clause is
added here rather than formalised from there, and it is the arithmetic of the
Mersenne weights that supplies it.

With $T_n=\sum_{k>n}\wgt k$, the standard
level-$n$ convex-hull cover consists of $2^n$ disjoint intervals of length
$T_n$, its nested intersection is $\Ach$, and
\[
  2^nT_n=\sum_{j\ge1}\frac{2^n}{2^{n+j}-1}\longrightarrow\sum_{j\ge1}2^{-j}=1,
\]
dominated by $2^{1-j}$.  Continuity of measure from above therefore gives
$\lambda(\Ach)=1$.

Membership is characterised level by level, by a greedy expansion.  Run through
$n=1,2,\dots$ carrying a remainder, initially the target $x$, and at level $n$
subtract $\wgt n$ from the remainder if $\wgt n$ does not exceed it, leaving
the remainder unchanged otherwise; call a level at which nothing is subtracted
\emph{skipped}.  Say that the expansion \emph{survives} level $n$ if the
remainder after that level is at most the remaining mass $T_n$.  Then
$x\in\Ach$ if and only if $x\ge0$ and the expansion survives every level
(\mword{Erdos249257/GreedyAchievementSet.lean}{1458}{mem_mersenneAchievementSet_iff_greedy_survival}{Lean}).

Non-membership of a nonnegative target is therefore visible at a single level,
and for a rational target strict failure can be certified effectively:
rational upper bounds for the remaining tail converge to $T_n$, so once one
lies below the rational residual the fatal inequality is proved.  The finite
example below uses an exact rational upper bound.  For
$x=3/4$, the greedy algorithm skips $\wgt1=1$ and leaves residual $3/4$,
while the exact zero-lookahead upper bound for the remaining tail is
$2\wgt2=2/3<3/4$; hence $3/4\notin\Ach$
(\mword{Erdos249257/GreedyAchievementSet.lean}{1761}{three_fourths_not_mem_mersenneAchievementSet}{Lean}).

The selector-to-subsum map sending a $0/1$ string $(\varepsilon_n)$ to
$\sum_n\varepsilon_n\wgt n$, which the strict tail inequality makes injective,
remains injective after restriction to any subfamily.  That matters because
Problem~\ref{res:problem} quantifies
over every infinite support rather than over $\Npos$.  For a set $J$ of future
offsets write $T_J(n)=\sum_{k\in J}\wgt{n+k+1}$ for the tail restricted to $J$
(\lword{Erdos257/MersenneSubseriesRigidity.lean}{20}{selectedMersenneTail}{definition},
summable at
\lword{Erdos257/MersenneSubseriesRigidity.lean}{24}{summable_selectedMersenneTail}{Lean}).

Then $T_J(n)<\wgt n$ for every $J$ and every $n\ge1$
(\lword{Erdos257/MersenneSubseriesRigidity.lean}{30}{selectedMersenneTail_lt_weight}{checked}),
since deleting weights only shrinks a tail that already sits below $\wgt n$ at
the full support.  Consequently the digit map is injective on strings
supported in $J$
(\lword{Erdos257/MersenneSubseriesRigidity.lean}{54}{supportedMersenneDigitValue_injective}{checked}),
stated on the subtype of
\lword{Erdos257/MersenneSubseriesRigidity.lean}{45}{SupportedMersenneDigits}{digit strings vanishing off $J$}
and evaluated by the
\lword{Erdos257/MersenneSubseriesRigidity.lean}{49}{supportedMersenneDigitValue}{restricted digit map},
so the statement is about the subseries itself and not a projection of the
full one.  Thus uniqueness of the selector survives passage to an arbitrary
subfamily; this is an injectivity statement, not a statement about binary
digits of the value.

That is a constraint on the shape of an argument, not on the supports: it
decides no value, and it is weaker than any statement in
Section~\ref{sec:forced}.

\subsection{Geometry after restricting the allowed exponents}

The injectivity statement has a geometric completion for every set
$J\subseteq\N$ of allowed digit positions.  Let
\[
  \Ach_J=
  \left\{\sum_{k\in J}\varepsilon_k\wgt{k+1}:
    \varepsilon_k\in\{0,1\}\right\}.
\]
Formally, the allowed strings form the
\lword{Erdos257/MersenneSubseriesRigidity.lean}{63}{supportedMersenneDigitSet}{supported digit set},
which is
\lword{Erdos257/MersenneSubseriesRigidity.lean}{66}{isClosed_supportedMersenneDigitSet}{closed}.
Its range is the
\lword{Erdos257/MersenneSubseriesRigidity.lean}{76}{supportedMersenneAchievementSet}{restricted achievement set};
the range and image descriptions agree by
\lword{Erdos257/MersenneSubseriesRigidity.lean}{79}{supportedMersenneAchievementSet_eq_image}{the image theorem}.
Consequently $\Ach_J$ is
\lword{Erdos257/MersenneSubseriesRigidity.lean}{90}{isCompact_supportedMersenneAchievementSet}{compact}
and
\lword{Erdos257/MersenneSubseriesRigidity.lean}{96}{isClosed_supportedMersenneAchievementSet}{closed}.

It lies inside $\Ach$ by
\lword{Erdos257/MersenneSubseriesRigidity.lean}{103}{supportedMersenneAchievementSet_subset}{support restriction}
and is therefore
\lword{Erdos257/MersenneSubseriesRigidity.lean}{112}{isNowhereDense_supportedMersenneAchievementSet}{nowhere dense}.
If $J$ is infinite, the digit space is
\lword{Erdos257/MersenneSubseriesRigidity.lean}{120}{preperfect_supportedMersenneDigitSet}{preperfect};
injectivity transfers that property to
\lword{Erdos257/MersenneSubseriesRigidity.lean}{150}{preperfect_supportedMersenneAchievementSet}{$\Ach_J$},
so the closed set is
\lword{Erdos257/MersenneSubseriesRigidity.lean}{167}{perfect_supportedMersenneAchievementSet}{perfect}.
Thus every infinite allowed support gives a compact perfect nowhere-dense set
with a unique selector for each point.

The metric classification is exact.  The summability lemma
\lword{Erdos257/MersenneSubseriesRigidity.lean}{175}{summable_mersenneDigitTerm}{controls digit terms},
and changing one coordinate changes the value by precisely its signed weight
by
\lword{Erdos257/MersenneSubseriesRigidity.lean}{182}{positiveMersenneDigitValue_update}{the update formula}.
If $k\notin J$, allowing $k$ splits the new set into $\Ach_J$ and its translate
by $\wgt{k+1}$
\lword{Erdos257/MersenneSubseriesRigidity.lean}{203}{supportedMersenneAchievementSet_insert}{exactly}.
The two pieces are
\lword{Erdos257/MersenneSubseriesRigidity.lean}{262}{disjoint_supportedMersenneAchievementSet_translate}{disjoint},
so adjoining one coordinate
\lword{Erdos257/MersenneSubseriesRigidity.lean}{286}{volume_supportedMersenneAchievementSet_insert}{doubles the volume}.
At $J=\N$, the restricted set is the full set $\Ach$
\lword{Erdos257/MersenneSubseriesRigidity.lean}{300}{supportedMersenneAchievementSet_univ}{by definition}.

\begin{theorem}[volume dichotomy for every allowed support]\label{res:supportvolume}
For every $J\subseteq\N$, exactly one of the following measure formulas
applies:
\[
\begin{array}{ll}
J=F^{\,c}\text{ for a finite }F,
  &\lambda(\Ach_J)=2^{-|F|},\\[2mm]
J^c\text{ is infinite},
  &\lambda(\Ach_J)=0.
\end{array}
\]
\end{theorem}

\noindent For finite $F$, the division-free identity
\lword{Erdos257/MersenneSubseriesRigidity.lean}{313}{volume_supportedMersenneAchievementSet_finset_compl_scaled}{multiplies the face volume by $2^{|F|}$}
and its solved form
\lword{Erdos257/MersenneSubseriesRigidity.lean}{349}{volume_supportedMersenneAchievementSet_finset_compl}{gives $2^{-|F|}$}.
Monotonicity under enlarging $J$ is
\lword{Erdos257/MersenneSubseriesRigidity.lean}{358}{supportedMersenneAchievementSet_mono}{checked}.
It traps a set with infinitely many forbidden coordinates below
finite-codimension faces of arbitrarily small dyadic measure, yielding
\lword{Erdos257/MersenneSubseriesRigidity.lean}{368}{volume_supportedMersenneAchievementSet_eq_zero_of_compl_infinite}{measure zero}.
The two cases are assembled in
\lword{Erdos257/MersenneSubseriesRigidity.lean}{397}{volume_supportedMersenneAchievementSet_dichotomy}{the formal dichotomy}.
This theorem classifies the size of the value set generated inside any
prescribed family of exponents.  It does not classify the arithmetic nature
of its individual points and therefore does not settle Problem~\ref{res:problem}.
An ordinary Hausdorff-dimension computation, recorded with
\texttt{ErdosProblems/Erdos257/RestrictedAchievementHausdorffDimension.md},
gives $\dim_H\Ach_J=\liminf_{N\to\infty}\#\{J\cap[1,N]\}/N$ for every
allowed index set $J$.  The arithmetic-progression case $J=m\Npos$ recovers
dimension $1/m$.  When $J$ is periodic of period $d$ with $0<r<d$ allowed
residues, the same note now records Ahlfors $r/d$-regularity and a biased
entropy dimension; this is a geometric strengthening, not an arithmetic
exclusion.  Fair-coin almost-sure irrationality on infinite $J$ remains
Kova\v{c}--Tao.

\subsection{The value \texorpdfstring{$1/2$}{1/2}}

A rational point of $\Ach$ attained by an infinite support refutes
Problem~\ref{res:problem}.  The distinguished candidate is $1/2$, and
Theorem~\ref{res:half} gives a self-contained pair of exact alternatives.

The greedy expansion of $1/2$ is an exact finite computation as far as one
cares to take it.  Through level $21$ it takes the exponents
\[
  2,\ 3,\ 6,\ 7,\ 14,\ 20,\ 21
\]
and skips the others, the skipped runs being $\{1\}$, $\{4,5\}$,
$\{8,\dots,13\}$ and $\{15,\dots,19\}$; every level through $21$ is survived.
The second condition below asks whether the list of skipped exponents is
infinite, and no finite computation answers that.

In the statement, $u=(u_1,\dots,u_d)\in\{0,1\}^{d}$ is a finite prefix,
$V(u)=\sum_{n\le d}u_n\wgt n$ is its value, and $T_n=\sum_{k>n}\wgt k$ as
above.

\begin{theorem}[membership and non-membership at $1/2$]\label{res:half}
The value $1/2$ belongs to $\Ach$ if and only if its canonical greedy
expansion omits infinitely many exponents.  Dually, $1/2\notin\Ach$ is
equivalent to the existence of a finite \emph{fatal gap}, a prefix $u$
through rank $d$ with
$V(u)+T_{d+1}<\tfrac12<V(u)+\wgt{d+1}$, which makes every continuation miss
and to the greedy orbit having a last skipped exponent.
Moreover no finite support has value $1/2$.
\end{theorem}

\noindent Checked as
\mword{Erdos249257/GreedyAchievementSet.lean}{2527}{half_mem_mersenneAchievementSet_iff_greedySkippedSupport_infinite}{the greedy form},
\mword{Erdos249257/HalfCylinderHalfMembershipClassification.lean}{126}{half_mem_mersenneAchievementSet_iff_unboundedTerminalFalse}{the terminal-bit form},
\mword{Erdos249257/HalfCylinderHalfMembershipClassification.lean}{213}{half_mem_mersenneAchievementSet_iff_exists_unboundedSkippedRanksAlong}{the skipped-rank form},
\mword{Erdos249257/HalfCylinderFatalGapRightTail.lean}{781}{seamGreedyEventuallyRight_iff_existsFatalHalfGap}{the fatal-gap equivalence},
\mword{Erdos249257/HalfCylinderFatalGapRightTail.lean}{787}{seamGreedyEventuallyRight_iff_half_not_mem}{its transfer to non-membership},
and
\mword{Erdos249257/HalfCarryReachability.lean}{589}{finite_boolSupport_ne_half}{the finite-support exclusion}.

\paragraph{Exact-row transport at the endpoint.}
The formal argument exposes the mechanism behind one conditional
half-membership implication.  Write $r_{c-1}$ for the rational greedy
remainder after ranks through $c-1$.  Whenever $c\ge4$ and

\[
  0<r_{c-1}<\wgt c,
\]

the skipped-core constructor first turns the finite greedy core into an exact
local quotient row at endpoint $2c-2$,
\mword{Erdos249257/BooleanMobiusSkipRowCofinal.lean}{55}{exactLocalMersenneHalfRow_of_positiveHalfGreedySkip}{the local row constructor}.
The hard finite step is explicit: the deficit below $1/2$ is below the next
Mersenne weight, so a Boolean word fills the entire upper-half capacity,
\mword{Erdos249257/BooleanMobiusSkippedCoreExactRow.lean}{228}{exists_upperHalfBooleanWord_of_skippedCore}{the upper-half Boolean fill}.
The residual cannot vanish at any finite rank,
\mword{Erdos249257/BooleanMobiusSkipRowCofinal.lean}{32}{greedyMersenneRemainderRat_half_pos}{the strict-positivity theorem};
the odd denominator of every finite positive Mersenne sum rules out exact
finite exhaustion of $1/2$.
If positive skips occur cofinally in the exact sense
\mword{Erdos249257/BooleanMobiusSkipRowCofinal.lean}{22}{CofinalPositiveHalfGreedySkips}{the cofinal-skip hypothesis},
their endpoints are cofinal by
\mword{Erdos249257/BooleanMobiusSkipRowCofinal.lean}{84}{cofinalExactLocalMersenneHalfRows_of_positiveHalfGreedySkips}{the row fan-in}.
The row-value estimate makes the discrepancy at endpoint $n$ at most
$(n+1)/2^n$, and closedness then places $1/2$ in $\Ach$,
\mword{Erdos249257/BooleanMobiusCofinalExactRows.lean}{71}{half_mem_mersenneAchievementSet_of_cofinalExactLocalRows}{the closed-set step}.

No compatibility between the finite witnesses is assumed.  The global
\mword{Erdos249257/BooleanMobiusSkipRowCofinal.lean}{97}{half_mem_mersenneAchievementSet_of_positiveHalfGreedySkips}{forward endpoint theorem}
is in fact reversible.

\begin{theorem}[positive skips characterise the half-value]
\label{res:positiveskipexact}
The condition \texttt{CofinalPositiveHalfGreedySkips} holds if and only if
$1/2\in\Ach$.
\end{theorem}

\noindent This is
\mword{Erdos249257/BooleanMobiusSkipRowCofinal.lean}{110}{cofinalPositiveHalfGreedySkips_iff_half_mem}{the checked equivalence}.
Thus cofinal positive skips are not an independently weaker hypothesis:
proving them already proves the half-membership endpoint, while
half-membership forces them because finite residuals never vanish.  Neither
side is established here, so the theorem identifies the exact obstruction
without supplying an unconditional counterexample to Problem~\#257.

\begin{theorem}[terminal scaled vanishing implies the half-value]
\label{res:terminalhalf}
Let $S$ be a
\texttt{HalfTerminalOnlyScaledVanishingSequence}: its finite words exclude
ranks $0$ and $1$, their depths tend to infinity, and the absolute terminal
carry divided by $2^M$ tends to zero along the depths $M$.  Then there is an
infinite set $A\subseteq\N$ with

\[
  \sum_{a\in A}\frac1{2^a-1}=\frac12.
\]
Consequently the universal irrationality assertion in
Problem~\ref{res:problem} is false under this hypothesis.  This is a
conditional implication, not a construction of $S$ and not a solution
of Problem~\#257.
\end{theorem}

\noindent The terminal-only argument first places $1/2$ in the achievement
set by the vanishing normalized-carry estimate, then uses the finite-support
exclusion to show that the limiting support is infinite:
\mword{Erdos249257/TerminalOnlyScaledVanishing.lean}{165}{half_mem_mersenneAchievementSet_of_terminalScaledVanishing}{achievement-set conclusion},
\mword{Erdos249257/TerminalOnlyScaledVanishing.lean}{221}{exists_infinite_support_half_of_terminalScaledVanishing}{infinite-support lift}.
The problem-level refutation is the direct formal consequence
\mword{Erdos257/HalfCounterexampleFrontier.lean}{39}{exists_rational_half_counterexample_of_terminalScaledVanishing}{rational half counterexample},
\mword{Erdos257/HalfCounterexampleFrontier.lean}{47}{not_universal_of_terminalScaledVanishing}{universal-claim refutation}.
The Mathlib-only Comparator statement packages those two problem-level
consequences in
\path{ExternalVerification257TerminalScaledVanishing/comparator.json}.
Its declaration is
\begin{center}
\small
\texttt{Erdos249257.ExternalVerification257TerminalScaledVanishing.}\allowbreak
\texttt{terminalScaledVanishing\_completeCounterexample},
\end{center}
and its reader-facing route is exactly Theorem~\ref{res:terminalhalf}.

For the mechanism-specific endpoint, write
\(R_s\) for the actual integer greedy remainder at seam rank \(s\).  After an
upper transition at \(d\), let \(p_{d,q}\) be the below-pulse at the cut of
rank \(d+q+1\), and define the exact affine charge by
\[
  C_{d,0}=0,\qquad
  C_{d,j+1}=4C_{d,j}+p_{d,j}+4.
\]
The assertion that the following inequalities hold is always restricted to
\(d\ge13\), \(k\le d\), an actual upper transition at \(d\), and its literal
right-run recurrence
\[
 R_{d+q+2}+2^{d+q+2}+p_{d,q}+4=4R_{d+q+1}
 \qquad(0\le q<k).
\]

\begin{theorem}[actual upper-successor counterexample endpoint]
\label{res:actual-upper-successor}
The universal complete-packet lower envelope
\[
  4^k\,2(d+k)\le R_{d+k+1}+C_{d,k}
\]
is equivalent to the universal pulse-free successor lower envelope
\[
  2^{d+1}-2^{d-k+1}+2(d+k)\le R_{d+1}.
\]
If these equivalent conditions hold, there is an infinite set
\(A\subseteq\N\) such that
\[
  \sum_{a\in A}\frac1{2^a-1}=\frac12,
\]
and therefore the universal irrationality assertion in
Problem~\ref{res:problem} is false.  No complete-packet or successor lower
envelope is constructed here, and the parent problem remains open.
\end{theorem}

\noindent The packet/successor equivalence is the exact checked statement
\mword{Erdos257/FatalCriticalDangerEndpointReduction.lean}{291}{seamActualUpperRightPacketLinearEscape_iff_successorLinearEscape}{in the source};
its direct half-membership consumer is
\mword{Erdos257/FatalCriticalDangerEndpointReduction.lean}{394}{half_mem_mersenneAchievementSet_of_actualUpperSuccessorLinearEscape}{checked}.
The problem-level propagation supplies the
\mword{Erdos257/ActualUpperSuccessorCounterexampleEndpoint.lean}{24}{infinite-support rational half witness}{checked witness},
\mword{Erdos257/ActualUpperSuccessorCounterexampleEndpoint.lean}{42}{universal-claim refutation}{checked refutation},
and
\mword{Erdos257/ActualUpperSuccessorCounterexampleEndpoint.lean}{59}{actualUpperSuccessorLinearEscape_completeCounterexample}{composite endpoint}.
The two selected Comparator declarations are recorded in
\path{ExternalVerification257ActualUpperSuccessor/comparator.json}; they route
exactly to this theorem and preserve the same open producer boundary.

The formal source also proves equivalent terminal-bit and unbounded
skipped-rank formulations in its internal half-cylinder seam coordinate; the
two middle source links above are those versions.  Their definitions are not
needed for the self-contained greedy and fatal-gap statement used here.
The two sides are asymmetric: non-membership is witnessed by a \emph{finite}
fatal gap and is therefore semi-decidable, whereas membership is an infinite
condition.  Computing further can only raise a lower bound on where a fatal
gap could occur; it cannot establish survival.

The next local obstruction is sharper than a numerical search.  Suppose a
take occurs at rank $b$, the next take at rank $c>b+1$, and put
$q=2^b-1$ and $m=2^{c-1}$.  If $R$ is the reciprocal of the residual just
before the first take, then the final intervening skip is dyadically unsafe
exactly when

\[
  \frac{q(m-1)}{q+m-1}<R<\frac{qm}{q+m}.
\]

The interval has the exact width
$q^2/((q+m)(q+m-1))$.  For a single intervening skip, $c=b+2$ and
$m=2q+2$, giving the two-thirds band

\[
  \frac{q(2q+1)}{3q+1}<R<\frac{2q(q+1)}{3q+2},
\]

whose width is below $1/9$.  Thus the band is an exact localization of one
greedy danger, not a claim that the actual orbit enters it.

\begin{theorem}[the sharp local two-adic obstruction]
\label{res:twothirdsband}
Let $p,D,q$ be positive odd integers.  If

\[
 q(2q+1)p<2D(3q+1),
 \qquad
 2D(3q+2)<2pq(q+1),
\]

then $p\ge7$.  In particular, the odd numerator classes $p=1,3,5$ cannot
occupy the cleared two-thirds band.  The bound is sharp for these local
hypotheses: odd triples with $p=7$ do occur in the band.
\end{theorem}

\noindent The hard step is genuinely $2$-adic.  The band defect
$\Delta=6D-2pq$ is positive and divisible by $4$ when $p,D,q$ are odd, so
$\Delta\ge4$; combining that gap with the band inequalities forces $p>6$.
The exact formal boundaries are the general localization
\mword{Erdos249257/HalfGreedyTwoThirdsBand.lean}{88}{postTakeUnsafeAt_iff_band}{general band localization},
the single-skip equivalence
\mword{Erdos249257/HalfGreedyTwoThirdsBand.lean}{127}{singletonUnsafe_iff_twoThirdsBand}{two-thirds band},
and the sharp numerator conclusion
\mword{Erdos249257/HalfGreedyTwoThirdsBand.lean}{231}{seven_le_of_intBand_odd}{odd numerator bound}.
Integral reciprocal states are excluded from the same band by
\mword{Erdos249257/HalfGreedyTwoThirdsBand.lean}{185}{not_twoThirdsBand_of_int}{integral safety}.
This is a one-step local no-go result: it neither proves that the half-greedy
orbit reaches or avoids the band nor settles half-membership.

Two further conditional arguments sharpen the boundary around this half-value
problem.  They are logically distinct from the terminal scaled-vanishing
question: one seeks irrationality from a structured support, while the other
would produce a rational counterexample from cofinal geometric stages.

\begin{theorem}[orthogonal-petal sunflower criterion]
\label{res:sunflower}
Let $A\subseteq\N$ support an
\texttt{OrthogonalPetalBouquet} and satisfy
\texttt{SunflowerForcedSlotTailSelection}.  Then

\[
  \operatorname{Irrational}\!\left(\sum_{a\in A}\frac{1}{2^a-1}\right).
\]
The bouquet hypothesis packages a finite exceptional/core part and residual
petals that are nontrivial, pairwise coprime, and reciprocal-summable.  The
remaining selector hypothesis is a uniform tail condition: at every positive
length $K$ it supplies a shifted block whose binary carry is divisible by
$2^K$ while its residual tail is bounded.  Neither hypothesis is constructed
here, so this is a conditional reduction rather than a solution of
Problem~\#257.
\end{theorem}

\noindent The hard mechanism is the persistent-reduced-modulus averaging
argument: the selector supplies carries, and the support-series theorem
turns that supply into irrationality.  The exact formal consumers are
\mword{Erdos249257/SupportSunflowerDichotomy.lean}{531}{sunflowerForcedCarrySupply_of_orthogonalPetalBouquet}{forced carry supply}
and
\mword{Erdos249257/SupportSunflowerDichotomy.lean}{540}{irrational_erdosSupportSeries_of_orthogonalPetalBouquet}{irrationality endpoint}.
The source definition of the selector makes the missing analytic obligation
explicit at
\mword{Erdos249257/SupportSunflowerDichotomy.lean}{406}{SunflowerForcedSlotTailSelection}{uniform tail selector}.
That selector is unnecessary for this support class.  Every orthogonal-petal
bouquet is already reciprocal-summable, so the flagship supplies all-base
irrationality with no forced-slot hypothesis:
\mword{Erdos249257/BouquetReciprocalIrrationality.lean}{19}{OrthogonalPetalBouquet.summable\_reciprocalSupportTerm}{reciprocal summability}
and
\mword{Erdos249257/BouquetReciprocalIrrationality.lean}{85}{irrational\_erdosSupportSeries\_of\_orthogonalPetalBouquet\_allBase}{all-base irrationality}.

\begin{theorem}[the composite-dilation defect]
\label{res:compositedefect}
Let $A\subseteq\N$ and let $a,x\in\Npos$ with $a\in A$.  Define the foreign-divisor
defect
\[
  \delta_A(a,x)=\#\{d\mid ax:d\in A,\ d\nmid x,\ d\ne a\}.
\]
Then the divisor incidence satisfies the exact identity
\[
  \scA_A(ax)=\scA_A(x)+\mathbf 1_{a\nmid x}+\delta_A(a,x).
\]
If every member of $A$ is prime, then $\delta_A(a,x)=0$, so the familiar
prime-support dilation formula is recovered.  More generally, if $A$ has an
\texttt{OrthogonalPetalBouquet} witness $h_B$, then for every $i\in\N$,
\[
  \delta_A(\operatorname{ray}_{h_B}(i),x)\le
  |h_B.\mathrm{exceptional}|+
  \scA_{\operatorname{range}(h_B.\mathrm{petal})}(x),
\]
where $\operatorname{ray}_{h_B}(i)=h_B.\mathrm{core}(i)h_B.\mathrm{petal}(i)$.
\end{theorem}

\noindent The first identity is the finite divisor partition checked at
\mword{Erdos249257/CompositeDilationDefect.lean}{30}{supportCoeff_mul_eq_add_defect}{exact dilation identity}.
The hard point is that a composite multiplier can create support divisors
which are neither the multiplier itself nor old divisors of $x$; they are
recorded by $\delta_A$ rather than silently discarded.  The defect vanishes
under the prime-support hypothesis by
\mword{Erdos249257/CompositeDilationDefect.lean}{103}{compositeDilationDefect_eq_zero_of_prime_support}{prime-support no-defect},
and the corresponding incidence formula is
\mword{Erdos249257/CompositeDilationDefect.lean}{119}{supportCoeff_mul_prime_support}{prime specialization}.
For a bouquet, every non-exceptional foreign divisor injects into a unique
petal divisor of $x$, yielding the displayed budget through
\mword{Erdos249257/CompositeDilationDefect.lean}{133}{mem_compositeDilationDefect_bouquet_imp_exceptional_or_petal_dvd}{foreign-divisor classification}
and
\mword{Erdos249257/CompositeDilationDefect.lean}{151}{compositeDilationDefect_le_exceptional_add_petalCoeff}{defect bound}.
The correction is concrete already for $A=\{2,6\}$: multiplying $x=1$ by
the composite support element $6$ creates the additional support divisor $2$,
so $\delta_A(6,1)=1$
(\mword{Erdos249257/CompositeDilationDefect.lean}{218}{compositeDilationDefect_two_six_fixture}{two-six witness}).

Thus a composite bouquet cannot be analysed by copying the prime formula
until this explicit foreign channel has been budgeted.  This theorem provides
that local budget only; it neither bounds defects along an infinite orbit nor
supplies the uniform tail selector required by Theorem~\ref{res:sunflower}.

\begin{theorem}[cofinal cylinders imply an infinite half-support]
\label{res:cylinderhalf}
Suppose that for every $N$ there are $M,K$ with
$\max\{N,1\}\le M$ and a nonempty
\texttt{CylinderStage}~$K~M$.  Then there is an infinite set
$A\subseteq\N$ with $0\notin A$ and

\[
  \sum_{a\in A}\frac{1}{2^a-1}=\frac12.
\]
Thus the universal irrationality assertion in Problem~\#257 is false under
this cofinal-stage hypothesis.  The formal proof does not construct the
cofinal family of full-cylinder stages.
\end{theorem}

\noindent The proof first converts each full-cylinder stage into a weaker
terminal-only strip, then invokes closedness of the half-achievement set
and the positive-support lift.  The exact steps are
\mword{Erdos249257/SuffixCylinderTerminalOnlyBridge.lean}{264}{cofinalTerminalOnlyStrip_of_cofinalCylinderStages}{terminal-only bridge},
\mword{Erdos249257/SuffixCylinderTerminalOnlyBridge.lean}{277}{exists_infinite_support_half_of_cofinalCylinderStages}{infinite half-support},
and
\mword{Erdos249257/SuffixCylinderTerminalOnlyBridge.lean}{287}{exists_infinite_positive_support_half_of_cofinalCylinderStages}{positive-support lift}.
This is a conditional implication, not evidence that such stages exist;
proving or refuting their cofinality is the remaining task for this argument.

One natural construction for $1/2$ is ruled out.  The Boolean support selected by the
negative values of the M\"obius function has value exactly $1/2$ plus the
positive M\"obius tail, hence at least $1/2+1/63$
(\mword{Erdos249257/MobiusSignSupportNoGo.lean}{111}{tsum_negativeMobius_eq_half_add_positiveMobiusTail}{identity},
\mword{Erdos249257/MobiusSignSupportNoGo.lean}{150}{half_add_one_div_sixty_three_le_tsum_negativeMobius}{bound}).
That closes the sign-truncation construction; it excludes no other support.

\subsection{A second rational target}

The target $1/21$ has a useful property that does not depend on any finite
search.

\begin{theorem}[finite-support obstruction at $1/21$]\label{res:one-over-twenty-one}
For every finite set $F$ of exponents with $n\ge2$ for all $n\in F$,
\[
  \sum_{n\in F}\frac{1}{2^n-1}\ne\frac1{21}.
\]
\end{theorem}

\noindent This is checked by
\lword{Erdos257/HalfCounterexampleFrontier.lean}{61}{finiteErdosSum_ne_one_div_twenty_one}{the formal theorem}.
The reduced denominator $21$ has doubling order $6$, so the exact
denominator--lcm identity forces every selected exponent to divide $6$.
The lower bound $n\ge2$ leaves only $2,3,6$, and the remaining eight subsets
are discharged by exact rational arithmetic.  Thus any representation of
$1/21$ along this rank-$\ge2$ route, if one exists, must be infinite.
The theorem proves nontermination only: it does not prove that $1/21$ lies
in $\Ach$.

A separate arithmetic lemma isolates the primitive cone that appears in one
candidate expansion.  Write
\[
  \mathcal P_{23}(n)=
  \{(p,q)\in\Npos^2:\gcd(p,q)=1,\ 2p+3q=n\}.
\]
Every $n\ge11$ has a member of $\mathcal P_{23}(n)$
(\mword{Erdos249257/Primitive23Multiplicity.lean}{52}{exists_primitive23_solution_of_eleven_le}{checked}),
whereas $\mathcal P_{23}(10)$ is empty
(\mword{Erdos249257/Primitive23Multiplicity.lean}{26}{no_primitive23_solution_ten}{checked}).
Multiplicity begins immediately at $11$, and every $10k$ with $k\ge2$ has
at least two distinct primitive representations
(\mword{Erdos249257/Primitive23Multiplicity.lean}{38}{exists_two_primitive23_solutions_eleven}{rank eleven},
\mword{Erdos249257/Primitive23Multiplicity.lean}{86}{exists_two_primitive23_solutions_mul_ten}{multiples of ten}).
These are exact Diophantine facts, not a counterexample construction.  In
particular the recurring collisions show why primitive-cone coverage is not
already a Boolean expansion: the missing step is a proof that the relevant
integer multiplicities can be converted into zero--one reciprocal-Mersenne
digits without changing the value.

\subsection{The scaled-greedy criterion at \texorpdfstring{$1/21$}{1/21}}

The finite-support obstruction makes $1/21$ useful only if the infinite
problem is stated in coordinates that retain the actual greedy orbit.  Let
\[
 r_N=\operatorname{greedyMersenneRemainder}(1/21,N)
\]
be the real remainder after rank $N$, put $y_N=2^Nr_N$, and set
\[
 c_{N+1}=\frac{2^{N+1}}{2^{N+1}-1}.
\]
The scaled recurrence is a nearly doubling map: it sends $y_N$ to
$2y_N-c_{N+1}$ when $c_{N+1}\le2y_N$, and to $2y_N$ otherwise.  The latter is
the exact lower branch, so rank $N+1$ is skipped precisely when
$2y_N<c_{N+1}$.

\begin{theorem}[scaled-greedy trap, escape, and cofinal return]
\label{res:scaled-greedy-trap}
For every nonnegative real $x$, writing $y_N(x)=2^N
\operatorname{greedyMersenneRemainder}(x,N)$,
\[
 x\in\Ach\quad\Longleftrightarrow\quad
 y_N(x)<2\quad\hbox{for every }N.
\]
If $x\notin\Ach$, then $y_N(x)\to+\infty$.  Conversely, a single bounded
cofinal subsequence already characterizes membership:
\[
 x\in\Ach\quad\Longleftrightarrow\quad
 \exists B\in\R\ \forall K\ \exists N\ge K:\qquad y_N(x)\le B.
\]
For every nonnegative rational $q$ this is also equivalent to
\[
 \forall K\ \exists N\ge K:\qquad
 2y_N(q)<\frac{2^{N+1}}{2^{N+1}-1}.
\]
In particular, $1/21\in\Ach$ if and only if its scaled greedy orbit crosses
this moving lower separatrix beyond every cutoff.
\end{theorem}

\noindent The trap-set equality is
\mword{Erdos249257/GreedyTrapDynamics.lean}{262}{mersenneAchievementSet_eq_scaledGreedyTrap}{kernel checked},
the escape theorem is
\mword{Erdos249257/GreedyTrapDynamics.lean}{189}{scaledGreedyRemainder_tendsto_atTop_of_not_mem}{kernel checked},
and the bounded-cofinal-return equivalence is
\mword{Erdos249257/GreedyTrapDynamics.lean}{225}{mem_mersenneAchievementSet_iff_scaledRemainder_cofinallyBounded}{kernel checked}.
The escape proof locates a fatal rank at which the residual exceeds the
complete remaining tail by some $\delta>0$.  That excess persists, so the
scaled residual is at least $2^N\delta$ thereafter.  A bounded cofinal
subsequence therefore forces membership.  Whether the $1/2$ or $1/21$
orbit has such a subsequence remains open.  Finite-horizon shadowing
witnesses may depend on the horizon, and the exact second-channel
recurrence retains the translations $2^{n+1}$; see
\texttt{ErdosProblems/Erdos257/HalfTargetQuantifierDiscipline.md}.

The general rational criterion is
\mword{Erdos249257/GreedyTrapDynamics.lean}{152}{rat_mem_mersenneAchievementSet_iff_scaledLowerBranchCofinally}{kernel checked},
and its $1/21$ specialization is
\mword{Erdos249257/GreedyTrapDynamics.lean}{283}{one_div_twentyOne_mem_iff_scaledLowerBranchCofinally}{kernel checked}.
The hard direction uses the rational normal form: membership is equivalent to
infinitely many omitted exponents, and a skip is exactly a lower-separatrix
crossing.  This is an exact reformulation, not a proof that the crossings
occur.  It asks for neither convergence nor a prescribed uniform bound.

The quotient coordinate gives a complementary classification of what
nonmembership would have to look like.  Let
\[
 Q_N=\left\lfloor\frac{2^N}{21}\right\rfloor-P_N,
\]
where $P_N$ is the integral numerator of the binary divisor-incidence prefix
generated by that same greedy support.  Thus $Q_N$ is the nonnegative
denominator-$21$ defect after the six-periodic residue of $2^N\bmod21$ has
been removed.  These are not independent models: the formal source proves an
exact identity relating $Q_N$, $2^Nr_N$ and a finite-prefix divisor tail
(\mword{Erdos249257/BooleanMobiusCarry.lean}{2117}{twentyOneGreedyDefect_add_mod_div_eq_scaled_remainder_add_tail}{checked}).

A finer denominator-specific coordinate works directly with the
denominator-$21$ quotient-greedy orbit.
At even depth $2R$, write
\[
 \begin{aligned}
 D_R&=\operatorname{twentyOneEvenQuotientGreedySupport}(R),\\
 s_R&=\operatorname{twentyOneEvenQuotientGreedyRemainder}(R).
 \end{aligned}
\]
Here $D_R\subseteq\{2,\ldots,R\}$ is the Boolean support obtained by descending
integer-greedy selection on the exact scaled quotient weights, and $s_R$ is its
terminal scalar remainder
(\mword{Erdos249257/TwentyOneQuotientGreedy.lean}{32}{twentyOneEvenQuotientGreedySupport}{support},
\mword{Erdos249257/TwentyOneQuotientGreedy.lean}{39}{twentyOneEvenQuotientGreedyRemainder}{remainder}).
Write $\mathcal F_{21}$ for \texttt{TwentyOneFatalAlignedBranch}, the explicit
branch in which the real greedy orbit has a fatal witness, only finitely many
skipped exponents (hence cofinite eventual selection), eventual
quotient/rational-greedy alignment, and eventual occupation of every doubling
block.

\begin{theorem}[fatal-branch quotient refinement at $1/21$]
\label{res:one-over-twenty-one-frontier}
The following statements hold.
\begin{enumerate}
\item $1/21\in\Ach$ if and only if $\mathcal F_{21}$ does not hold.
\item If there is an unbounded sequence of ranks $R$ with $s_R\le 2^R$,
  then $1/21\in\Ach$.
\item On $\mathcal F_{21}$, eventually $s_R>2^R$, the boundary rank $R+1$
  belongs to $D_{R+1}$, and $(D_R,s_R)$ follows one exact affine recurrence.
\end{enumerate}
\end{theorem}

\noindent The equivalence is
\mword{Erdos249257/TwentyOneQuotientGreedy.lean}{3507}{one_div_twenty_one_mem_iff_not_fatalAlignedBranch}{kernel checked};
the closed-row compactness step is
\mword{Erdos249257/TwentyOneQuotientGreedy.lean}{5554}{twentyOneCofinalEvenQuotientGreedyDecay_of_closedRows}{kernel checked};
and the eventual affine regime is
\mword{Erdos249257/TwentyOneQuotientGreedy.lean}{5658}{twentyOneFatalAlignedBranch_eventually_affine_supercapacity}{kernel checked}.
Explicitly, on $\mathcal F_{21}$ the quotient orbit is eventually strictly
supercapacity and loses its last Boolean choice: if $\tau_R$ is the periodic
target pulse and $\pi_R$ the divisor pulse generated by $D_R$, then for every
sufficiently large $R$,
\[
 s_R>2^R,\qquad D_{R+1}=D_R\cup\{R+1\},\qquad
 s_{R+1}=4s_R+\tau_R-\pi_R-(2^{R+1}+1).
\]
The denominator-specific separation theorem additionally proves that every
closed Boolean quotient row is exactly the canonical quotient-greedy row,
including exact saturation at the boundary
(\mword{Erdos249257/TwentyOneQuotientGreedy.lean}{231}{twentyOneClosedRow_forces_quotientGreedy}{kernel checked}).
An aligned crossing from saturation into strict supercapacity forces a
missing canonical ancestor at one-third or two-thirds scale and a real
skipped exponent with square-root-bounded defect
(\mword{Erdos249257/TwentyOneQuotientGreedy.lean}{5179}{twentyOneAlignedSaturatedCrossing_forces_canonical_ancestor_hole}{ancestor hole},
\mword{Erdos249257/TwentyOneQuotientGreedy.lean}{5223}{twentyOneAlignedSaturatedCrossing_forces_scaled_greedy_skip}{scaled skip}).

These conclusions remove alternative late Boolean branches; they do not
exclude the full fatal/cofinite/aligned branch.  On that branch the permanent
affine-supercapacity recurrence is forced, so a contradiction of the
recurrence would be sufficient, but the recurrence alone is not the exact
membership endpoint.  It is conditional, not a contradiction: synthetic pulses
can sustain such an affine recurrence without being the actual divisor pulse.

\section{Open problems}\label{sec:open}

Problem~\ref{res:problem} remains the frame: the settled families in
Section~\ref{sec:map} do not approach a universal quantifier, and the forced
conditions in Section~\ref{sec:forced} are not known to be contradictory.
Theorem~\ref{res:reciprocal-support} does, however, force every counterexample
into the reciprocal-divergent region
$\sum_{a\in A}a^{-1}=\infty$.
For the base-$2$ universal problem, three endpoint questions organise the
remaining discussion.  The universal problem and the half-value question
are not independent, and the relation between them is asymmetric.  No finite
support has value $1/2$
(Theorem~\ref{res:half}), so $1/2\in\Ach$ would exhibit an \emph{infinite}
support with a rational value and thereby refute
Problem~\ref{res:problem}; equivalently, a positive answer to
Problem~\ref{res:problem} puts $1/2$ outside $\Ach$.  The converse direction
does not follow: $1/2\notin\Ach$ would give a finite fatal-gap witness and
eliminate this candidate, but would not imply the universal statement.  The
half-value question is therefore a one-sided test of \#257, not a second
problem beside it.

\begin{problem}[membership of 1/21 in the Mersenne achievement set]
\label{prob:one-over-twenty-one-membership}
Prove cofinal crossings of the exact moving lower separatrix
\[
 2\,\operatorname{scaledGreedyRemainder}(1/21,N)
 <\operatorname{mersenneScale}(N+1).
\]
Equivalently, exclude the explicit fatal/cofinite/aligned branch $\mathcal F_{21}$.
It is sufficient either to contradict the eventual permanent
affine-supercapacity recurrence forced by that branch or to force an unbounded
sequence of closed canonical quotient rows; neither sufficient route is
claimed to be equivalent by itself.
\end{problem}

\begin{enumerate}
\item \textbf{Problem~\ref{res:problem} itself}, for arbitrary infinite $A$.
  Theorem~\ref{res:reciprocal-support} excludes every reciprocal-summable
  support, so an unresolved counterexample must have divergent reciprocal
  mass.  Section~\ref{sec:forced} constrains every hypothetical counterexample
  without excluding one in that remaining region.  Every counterexample would have
  unbounded scaled-tail states, sublogarithmic divisor-coverage gaps and an
  admissible Boolean--M\"obius scaled-tail sequence.  The conditional lower
  bounds in Theorem~\ref{res:mass} are now subordinate inside the proved
  reciprocal-summable exclusion rather than the frontier itself.  No
  contradiction is known in the reciprocal-divergent region.
\item \textbf{Membership of $1/2$ in $\Ach$.}  Theorem~\ref{res:half} makes
  this exactly equivalent to infinitely many greedy skips, and to a fatal gap
  on the other side; neither is proved.  A proof of non-membership would take
  the form of one finite fatal gap.
\item \textbf{Problem~\ref{prob:one-over-twenty-one-membership}.}
  Theorem~\ref{res:one-over-twenty-one} rules out finite support, while
  Theorem~\ref{res:one-over-twenty-one-frontier} reduces membership exactly
  to excluding $\mathcal F_{21}$.  Contradicting its forced eventual
  affine-supercapacity recurrence or producing arbitrarily deep closed
  canonical rows would suffice.  Neither input is known, and neither is
  promoted to an equivalence.
\end{enumerate}

\noindent The half-value question remains valid on both sides; the criterion at
$1/21$ is at present the sharper of the two.  The questions below refine that
criterion and the arithmetic constraints around it, in order of how directly
an answer would move the proved boundary.

\begin{problem}[permanent affine supercapacity]\label{prob:affine-supercapacity}
Can the \emph{actual} denominator-$21$ greedy orbit satisfy
$\mathcal F_{21}$ indefinitely?  Equivalently, can the complete fatal,
cofinite-selection, alignment and doubling-block hypotheses coexist with the
eventual recurrence
\[
 s_R>2^R,\quad D_{R+1}=D_R\cup\{R+1\},\quad
 s_{R+1}=4s_R+\tau_R-\pi_R-(2^{R+1}+1),
\]
or must an arbitrarily deep closed return $s_R\le2^R$ occur?
\end{problem}

\noindent A contradiction from a Lyapunov function, a $2$-adic obstruction,
a divisibility theorem or a recurrence classification would prove
$1/21\in\Ach$ and, by Theorem~\ref{res:one-over-twenty-one}, produce an
infinite rational support.  Conversely, an actual construction satisfying
\emph{all} clauses of $\mathcal F_{21}$ would prove $1/21\notin\Ach$.
Constructing only an abstract affine orbit with adversarial pulses does
neither.

\begin{problem}[weakest native recurrence criterion]\label{prob:scaled-return}
Does the scaled actual greedy remainder return cofinally to one bounded
interval?
\[
 \exists B<\infty\ \forall K\ \exists N\ge K:
 \qquad 2^N r_N\le B.
\]
\end{problem}

\noindent This is not merely sufficient: it is the $x=1/21$ specialization
of the general membership equivalence in
Theorem~\ref{res:scaled-greedy-trap}
(\mword{Erdos249257/GreedyTrapDynamics.lean}{275}{one_div_twentyOne_mem_iff_scaledRemainder_cofinallyBounded}{checked}).
It asks for neither convergence, a prescribed $B$, a global bound nor bounded
return gaps.  Two concrete stronger targets are cofinal recurrence of
$Q_N\le1$
(\mword{Erdos249257/BooleanMobiusCarry.lean}{2843}{one_div_twenty_one_mem_mersenneAchievementSet_of_defect_le_one_cofinally}{sufficient condition})
and arbitrarily deep closed quotient rows $s_R\le2^R$.  Exact computation
through rank $200{,}000$ records $96$ zero-defect returns and $4{,}956$ returns
to $Q_N\le1$, with last observed ranks $193{,}690$ and $199{,}930$ and maximum
observed gap $492$; these finite data do not prove cofinality.

\begin{problem}[actual-orbit invariant]\label{prob:actual-invariant}
Is there a finite-memory, $2$-adic or discrepancy invariant, using the
correlated divisor pulses of the actual support, that forces a closed return
or forbids permanent supercapacity?  More precisely, can one use a bounded
window of $R\bmod6$, residues of $s_R$, endpoint divisor counts and the finite
set of eventual skips to force descent or a forbidden state?  Alternatively,
can one prove that no bounded-memory invariant distinguishes the true orbit
from synthetic permanent-supercapacity controls?
\end{problem}

\noindent Pointwise pulse bounds are known to be too weak.  The exact
six-step recurrence is
\[
 Q_{N+6}=64Q_N+3(2^N\bmod21)-L_N,
\]
with $L_N$ the actual weighted repair load
(\mword{Erdos249257/BooleanMobiusCarry.lean}{1892}{twentyOneGreedyDefect_add_six}{checked}).
The formerly plausible translation-invariant six-step contraction first
fails at $N=73$; the surviving statement is the slope-aware repair-load iff
at
\mword{Erdos249257/BooleanMobiusCarry.lean}{1927}{twentyOneGreedyDefect_add_six_le_div_six_iff_repairLoad}{the exact equivalence}.

\begin{problem}[final-skip Diophantine exclusion]\label{prob:fatal-interval}
Let $E=\sum_{n\ge1}(2^n-1)^{-1}$.  If $1/21\notin\Ach$, let $M$ be the last
skipped exponent, $S_M$ its finite skipped prefix, and
\[
 a_M=\frac1{21}+\sum_{d\in S_M}\frac1{2^d-1}.
\]
Can the complete final-skip signatures be used to prove
$|E-a_M|\ge\operatorname{gap}_M$, contradicting
\[
 0<a_M-E<\operatorname{gap}_M?
\]
\end{problem}

\noindent The fatal interval is checked as
\mword{Erdos249257/TwentyOneQuotientGreedy.lean}{3735}{lastTwentyOneSkip_erdosBorwein_fatalInterval}{an exact one-sided approximation}.
The reduced skipped-prefix denominator has binary order equal to the lcm of
the skipped exponents and hence at least $M$
(\mword{Erdos249257/TwentyOneQuotientGreedy.lean}{3569}{lastTwentyOneSkip_oddDoublingOrder_eq_lcm}{order identity},
\mword{Erdos249257/TwentyOneQuotientGreedy.lean}{3583}{lastTwentyOneSkip_le_oddDoublingOrder}{lower bound}).
Parity-sensitive gcd restrictions and adjacent-numerator product inequalities
are also checked, but they are necessary signatures rather than the missing
upper-height or irrationality-measure estimate.

\begin{problem}[classification of rational targets]\label{prob:target-classification}
For $L\ge1$, classify the reduced rationals
\[
 \mathcal T_L=\left\{\frac pq\in(0,1):
 \begin{array}{l}
 \gcd(p,q)=1,\ q\ \hbox{odd},\ \operatorname{ord}_q(2)\le L,\\
 p/q\ \hbox{has no finite Mersenne representation}
 \end{array}\right\}
\]
for which non-membership reduces to an eventually deterministic finite-memory
or affine quotient recurrence.  In particular, classify $\mathcal T_6$ by
target-pulse alphabet, finite-support obstruction, weakest membership
criterion, surviving branch complexity and available Diophantine estimate.
Is $1/21$ minimal or unique under explicit criteria, or does another target
have a strictly simpler fatal branch?
\end{problem}

\begin{problem}[realised finite denominators]\label{prob:denominator-realisation}
For fixed $b\ge2$ and $L\ge2$, let $\mathcal D_b(L)$ be the reduced
denominators of the nonempty finite sums
$\sum_{n\in F}(b^n-1)^{-1}$ with $\operatorname{lcm}(F)=L$.  Is
\[
 \mathcal D_b(L)=
 \{D:D\mid b^L-1,\ \operatorname{ord}_D(b)=L\}?
\]
If not, what cyclotomic, valuation or cancellation conditions characterise
the realised denominators?
\end{problem}

\noindent The proposed equality is false: $\mathcal D_2(6)=\{21,63\}$ even
though $\operatorname{ord}_9(2)=6$.  For prime period $p$ one has the complete
formula $\mathcal D_b(p)=\{b^p-1,(b^p-1)/\gcd(b-1,p+1)\}$.  These ordinary
facts are recorded with
\texttt{ErdosProblems/Erdos257/FiniteDenominatorRealisation.md}.  They do not
supply a height estimate for the actual final-skip argument.

\noindent For scope, all squarefree values at power-of-two bases, jointly in
each finite family, are settled by Corollary~1.2 and Example~1.1 of
~\cite[author-preprint p.~4]{duverneytachiya}
(Corollary~\ref{res:sfsettled}).  Their cited
corollary does not cover bases that are not powers of \(2\), and this note
makes no global open-status or priority claim for those values.

\section*{Statements and declarations}

This manuscript is authored exposition, not proof authority.  The linked Lean
snapshot is authoritative only for its exact propositions; kernel checking
establishes that a proposition was proved, not that it is interesting, novel,
or sufficient.  The displayed proof of Corollary~\ref{res:blind}, the
derivation of Corollary~\ref{res:sfsettled} from
~\cite[Cor.~1.2 and Ex.~1.1, author-preprint p.~4]{duverneytachiya},
and the general finite-change argument (from the two checked prefix lemmas) are
expository arguments rather than named checked statements.

The terminal scaled-vanishing Comparator package selects the single composite
declaration
\begin{center}
\small
\texttt{Erdos249257.ExternalVerification257TerminalScaledVanishing.}\allowbreak
\texttt{terminalScaledVanishing\_completeCounterexample}
\end{center}
and routes it to Theorem~\ref{res:terminalhalf}.  Its Challenge imports only
Mathlib; its Solution transports the exact infinite-support half-value and
universal-refutation endpoints named at that theorem.

The actual upper-successor Comparator package selects the exact equivalence
\begin{center}
\small
\texttt{Erdos249257.ExternalVerification257ActualUpperSuccessor.}\allowbreak
\texttt{actualUpperRightPacketLinearEscape\_iff\_successorLinearEscape}
\end{center}
and the composite endpoint
\begin{center}
\small
\texttt{Erdos249257.ExternalVerification257ActualUpperSuccessor.}\allowbreak
\texttt{actualUpperSuccessorLinearEscape\_completeCounterexample}.
\end{center}
Both route to Theorem~\ref{res:actual-upper-successor}.  The Mathlib-only
Challenge states the concrete integer greedy dynamics; the Solution transports
the exact source declarations without adding hypotheses or weakening either
conclusion.

The squarefree
no-go interfaces, shifted coefficient, shift equivalence, Chinese-remainder
block supply, and restricted-selector injectivity are directly checked
statements linked at their use; they are not prose-only claims.
Ordinary arguments recorded with
\texttt{ErdosProblems/Erdos257/VariableExponentCoverSeparation.md},
\texttt{ErdosProblems/Erdos257/StrengthenedVariableExponentCover.md},
\texttt{ErdosProblems/Erdos257/RestrictedAchievementHausdorffDimension.md},
\texttt{ErdosProblems/Erdos257/ElementarySupportConstraints.md},
\texttt{ErdosProblems/Erdos257/FiniteDenominatorRealisation.md},
\texttt{ErdosProblems/Erdos257/SubpowerRepairWindows.md},
\texttt{ErdosProblems/Erdos257/CoverFirstLogarithmicMoment.md}
and
\texttt{ErdosProblems/Erdos257/HalfTargetQuantifierDiscipline.md}
are reviewed notes, not Lean authority, and they do not replace
Theorem~\ref{res:reciprocal-support}.  The signed finite-period theorems are
Lean-checked in a focused module, not in the Palomar flagship.

The numerical
instances, the
table of finite supports in Section~\ref{sec:period}, the incidence and
zero-window examples, the reciprocal-mass and Boolean--M\"obius instances, the
$3/4$ certificate, and the greedy expansion of $1/2$ through level~$21$,
are direct computations, reported as computations and not as theorems.
Results attributed to Campbell, to Duverney and Tachiya, to Tao and
Ter\"av\"ainen, to Luca and Tachiya, and to Kova\v{c} and Tao are cited from
the literature and are not formalised here.

Erd\H{o}s Problem~\#257 remains open.

\begin{thebibliography}{10}
\bibitem{kakeya1914} S.~Kakeya,
  \href{https://doi.org/10.11429/ptmps1907.7.14_250}
  {\emph{On the set of partial sums of an infinite series}},
  Proc. Tokyo Math.-Phys. Soc., 2nd ser., \textbf{7} (1914), 250--251.
  The perfectness theorem and the real tail criterion are on p.~251.
\bibitem{duverneysuzukitachiya2019} D.~Duverney, Y.~Suzuki and Y.~Tachiya,
  \href{https://arxiv.org/abs/1908.07290}
  {\emph{Linear independence results for certain sums of reciprocals of
  Fibonacci and Lucas numbers}}, arXiv:1908.07290, 2019.
  Corollary~1 gives the stated sparse reciprocal-sum linear independence
  under the pairwise-coprime, odd-support and reciprocal-summability
  hypotheses.
\bibitem{campbell2026} J.~M. Campbell,
  \href{https://arxiv.org/abs/2605.24160}
  {\emph{On the binary digits of the Erd\H{o}s--Borwein constant}},
  arXiv:2605.24160v1, 2026.  Theorem~1 on p.~12 proves that the binary block
  \(11\) occurs infinitely often in the base-\(2\) expansion of the
  full-support value; its proof occupies pp.~12--24.
\bibitem{duverneytachiya} D.~Duverney and Y.~Tachiya,
  \href{https://danielduverney.fr/documents/theorie-des-nombres/DuverneyTachiya190522.pdf}
  {\emph{Refinement of the Chowla--Erd\H{o}s method and linear independence
  of certain Lambert series}},
  Forum Math.\ 31 (2019), no.~6, 1557--1566,
  \href{https://doi.org/10.1515/forum-2018-0299}{DOI}.  Corollary 1.2 gives
  the general $F_s(E)$ theorem and its monomial images under
  $|q|^{\lcm(1,\dots,\ell)}\le s$; Example 1.1 gives the joint squarefree
  family at all bases \(2^j\), \(j\ge1\).  Both are on p.~4 of the linked
  author preprint; the proof is on pp.~9--11.
\bibitem{taoteravainen2025} T.~Tao and J.~Ter\"av\"ainen,
  \emph{Quantitative correlations and some problems on prime factors of
  consecutive integers}, arXiv:2512.01739 (submitted December 2025, revised
  April 2026).  Theorem 1.3 proves $\sum_{n\ge1}\omega(n)/2^n=\sum_p(2^p-1)^{-1}$
  irrational, settling the prime-support case of \#257 at base~$2$; the
  extension to every integer base and the prime-power support are stated as
  remarks with the modifications left to the reader.  The theorem is on
  p.~4 and its proof is Section~5, pp.~44--56, in arXiv v2.
\bibitem{lucatachiya2017} F.~Luca and Y.~Tachiya,
  \href{https://www.kurims.kyoto-u.ac.jp/~kyodo/kokyuroku/contents/pdf/2014-14.pdf}
  {\emph{Linear independence results for the values of divisor functions
  series}}, RIMS K\^oky\^uroku No.~2014 (2017), 138--150.
  Theorem~A on p.~139 explicitly restates their earlier Theorem~1.1 for
  nonzero purely periodic integer weights.  Theorem~1 is on p.~139, its
  examples are on p.~140, and its proof is on pp.~149--150.
\bibitem{kovactao} V.~Kova\v{c} and T.~Tao, \emph{On several irrationality
  problems for Ahmes series}, Acta Math.\ Hungar.\ 175 (2025), 572--608,
  \href{https://doi.org/10.1007/s10474-025-01528-0}{DOI}.  Remark 4.1
  (p.~13) records the strict tail inequality and the Cantor structure.
  Theorem~2.3 (p.~5; proof pp.~13--14) constructs rational merged sums from
  several bases under its mass hypothesis; it is not a fixed-base
  counterexample.
\bibitem{erdos1948} P.~Erd\H{o}s, \emph{On arithmetical properties of Lambert
  series}, J.~Indian Math.\ Soc.\ 12 (1948), 63--66.
\bibitem{erdos1968} P.~Erd\H{o}s,
  \href{https://users.renyi.hu/~p_erdos/1969-09.pdf}{\emph{On the
  irrationality of certain series}}, Math.\ Student 36 (1968), 222--226
  (issued 1969).  The theorem on p.~222 treats pairwise-coprime support with
  convergent reciprocal sum at every integer base $b\ge2$; the claimed
  removal of pairwise coprimality is stated without proof.
\bibitem{erdosgraham1980} P.~Erd\H{o}s and R.~L. Graham,
  \href{https://mathweb.ucsd.edu/~ronspubs/80_11_number_theory.pdf}{\emph{Old
  and New Problems and Results in Combinatorial Number Theory}},
  Monogr. Enseign. Math. 28, Geneva, 1980, p.~62.
\bibitem{erdos1988} P.~Erd\H{o}s, \emph{On the irrationality of certain
  series: problems and results}, in A.~Baker (ed.), \emph{New Advances in
  Transcendence Theory}, Cambridge UP, 1988, pp.~102--109,
  doi:\href{https://doi.org/10.1017/CBO9780511897184.009}
  {10.1017/CBO9780511897184.009}.
\bibitem{erdosproblems} T.~F. Bloom,
  \href{https://www.erdosproblems.com/257}{\emph{Erd\H{o}s Problem \#257}},
  \texttt{erdosproblems.com/257}, accessed 28 July 2026 (page displays
  ``last edited 15 April 2026'').  The current record labels the universal
  fixed-base problem open, cites \texttt{[Er68d]},
  \texttt{[ErGr80, p.~62]} and \texttt{[Er88c, p.~105]}, records the
  Tao--Ter\"av\"ainen prime and prime-power cases and the Kova\v{c}--Tao
  refutation of a broader varying-denominator speculation, and explicitly
  describes its status as the website owner's present assessment rather than
  a literature-completeness guarantee.
\end{thebibliography}

\end{document}
